\maketitle
\noteprovenance

\begin{abstract}
\noindent
Erd\H{o}s, Herzog, and Piranian ask whether every monic polynomial of degree
$n\ge2$ with zeros in the open unit disc has two zeros joined inside
$\{|f|<1\}$ by a curve of length less than $2$.  Every monic trinomial
$f(z)=z^n+az^m+b$ with $1\le m<n$ and every zero in the open unit disc has the
entire segment from the origin to each zero inside $\{|f|<1\}$, so any two
zeros are joined through the origin by a broken line of length below $2$, in
every degree and with the middle coefficient unrestricted.  Separation of one
critical value gives the same conclusion for an arbitrary monic polynomial: if
$c$ is a simple critical point, $v=f(c)\ne0$, and $|1-f(d)/v|\ge2$ at every
other critical point, then two zeros are joined inside $\{|f|<1\}$ by a curve
of length below $2$ in every degree $n\ge3$.

The separation theorem resolves the saddle by a square-root coordinate.
Separation makes the resolved inverse branch $Z$ univalent on the disc of
radius $\sqrt S$, the identity $P(Z(\xi))=1-\xi^2$ pins the exact Dirichlet
integral $\int|P'|^2$ over its image, and Crane's sharp lower bound for the
area of a polynomial image then bounds the Taylor coefficients of $Z$ and
hence its length by
\[
 2\Bigl(\frac2n\Bigr)^{1/(2n)}|v|^{1/n}S^{1/n}
   \sqrt{\log\!\frac{S}{S-1}}.
\]
The threshold $S=2$ therefore reaches every degree $n\ge3$, where the earlier
form of the estimate reached only $n\ge6$.

Write $\mu=\min_{f'(c)=0}|f(c)|$ and $\rho=\mu^{1/n}$.  Two further statements
hold for arbitrary monic polynomials.  Every squarefree $f$ of degree $n\ge2$
with $\mu\le13/25$ has two distinct zeros joined inside $\{|f|<1\}$ by a
rectifiable curve of length strictly below $2$, with no hypothesis on where the
zeros lie; that proof is ordinary analysis closed by an exact rational
certificate with certified quantity $X=635762889599/1000000000000$ and closing
inequality $(13/25)e^{X}<1$.  With the threshold removed, two zero occurrences
of any monic $f$ are joined inside $\{|f|\le2\mu\}$ by a path of length at most
$(71/10)\rho$, and inside $\{|f|<1\}$ with length at most $5.7$ once the zeros
lie in the open unit disc and $\mu\le1/2$.  The same construction settles the
problem outright for $\mu\le1/2$ whenever the first-merge component carries at
least $17$ zeros, and whenever its capacity defect satisfies
$\operatorname{cap}(\overline C)\le(2\mu)^{1/n}/3$.

Degree three is settled in full.  We also prove the sharp adjacent-zero
envelope for collinear zeros, with affine Chebyshev equality configurations,
and solve the primitive sparse quintic and translated cubic quotient-fibre
families.  An auxiliary sharp power-sum bound
$\sum_{j=1}^{n-1}|f(c_j)|^{1/(n-1)}\le(n-1)R^{n/(n-1)}$ holds in every degree
$n\ge2$ for zeros in a closed disc of radius $R\ge0$.

The unrestricted problem remains open.  Critical values near a tie and multiple
saddles escape the separation argument, and the regime $13/25<\mu<1$ is
untouched by the threshold argument.  Exact counterexamples show that short
Euclidean chords and critical spokes need not stay in the lemniscate.  The
remaining task is to select and control a contained connector; the
critical-value budget alone does not bound its length.
\end{abstract}

\begin{paperresult}{Solved families and two all-degree regimes}
\textbf{Trinomials.} Theorem~\ref{res:trinomial-all-degree} settles the
Erd\H{o}s--Herzog--Piranian conclusion for every monic trinomial in every
degree, with a prescribed path and an unrestricted middle coefficient.  Its
inequalities are checked by the Lean kernel.
\textbf{Separated critical values.}
Theorem~\ref{res:critical-value-separation} bounds the length of the resolved
inverse-ray connector at a simple isolated critical value, and
Corollary~\ref{res:critical-value-thresholds} turns that bound into the target
conclusion at separation $S=2$ in every degree $n\ge3$.
\textbf{Regimes with no root hypothesis.}
Theorem~\ref{res:low-critical-thirteen-twentyfifths} settles every squarefree
monic polynomial with least critical-value modulus at most $13/25$, in every
degree.  Theorem~\ref{res:constant-factor-path} removes the threshold and pays
a constant factor, and its arity and capacity corollaries settle the problem
outright on two further explicit regions.
Theorem~\ref{res:degree-three} settles degree three.
\textbf{Complementary results.} Sharp solved families and the critical-value
budget through degree five supply geometric and algebraic information beyond
those regimes.
\textbf{Open boundary.} Near-tied critical values and multiple saddles
require another construction.  An admissible-hub selector is one sufficient
route to the target; no characterization of every possible solution is proved.
\end{paperresult}

\externalreviewobject{The paper proves the Erd\H{o}s--Herzog--Piranian
conclusion for every monic trinomial in every degree, with the radial
inequalities kernel-checked; a quantitative connector at a separated simple
critical value, whose threshold $S=2$ reaches every degree $n\ge3$ through
P\'olya's area inequality and Crane's sharp polynomial-image bound; a short
contained connector whenever the least critical-value modulus is at most
$13/25$; a constant-factor connector of length at most $(71/10)\mu^{1/n}$
inside $\{|f|\le2\mu\}$ with no hypothesis at all, together with arity and
capacity corollaries that reach the target constant; the full degree-three
case; a sharp all-degree collinear theorem; a complete primitive sparse
quintic theorem; and translated cubic quotient-fibre theorems in every degree
$3q$.  It also checks exponential Newton-value decay, positive-ray collision
interfaces, finite translation avoidance, and quantified root retention, and
records the current near-Fekete residual.}{The separation regime, the two
threshold-free theorems, the degree-three theorem and the three solved-family
assemblies are ordinary mathematics without kernel-checked authority.  The
$13/25$ threshold, the $(71/10)$ constant and the degree-three theorem have no
formal endpoint in the pinned corpus, and the regime $13/25<\mu<1$ is
untouched by the threshold argument.  The separation argument excludes
near-tied critical values and multiple saddles.  The checked dynamical and
perturbative inputs do not repair the global topology and metric gluing in the
complementary strata, so Problem~\#1041 remains open.  The degree-five target
and the no-go witnesses are boundaries around the open theorem and close
nothing.}{Comparator selects the trinomial theorem, the finite-family
small-translation theorem, the quantified root-retention theorem and three
solved-family kernels.  Critical-value separation, including its threshold
analysis, has no corresponding formal endpoint in the pinned corpus.  Each
formal endpoint routes to the exact paper result and boundary it supports; the
covering-space and area argument and the frontier section remain ordinary
evidence classes.}

The two selected perturbation interfaces have the following exact boundaries.
For a finite type $\iota$, an injective family of critical values
$c\colon\iota\to\mathbb C$, and every $\varepsilon>0$, the row
\texttt{exists\_small\_translation\_separating\_arguments} supplies a shift
of norm less than $\varepsilon$ such that every translated value is nonzero
and every two distinct translated values lie on different positive rays.
Separately, for a monic split polynomial of positive degree whose roots have
norm at most $\rho$, the row
\texttt{constant\_perturbation\_roots\_in\_unitDisk} requires $\rho\ge0$,
$\varepsilon>0$, the margin
\[
 ((\operatorname{natDegree}f+1)\varepsilon)^{1/\operatorname{natDegree}f}
   +\rho<1,
\]
and a shift of norm less than $\varepsilon$; it then places every root of
$f+\operatorname{C}(\text{shift})$ in the open unit disc.  These are finite
perturbation and stability interfaces only: neither row supplies the missing
global topology or the length-$2$ gluing argument.

\section{The problem}\label{sec:problem}

\begin{problem}[Erd\H{o}s \#1041]\label{res:problem}
Let $f(z)=\prod_{i=1}^{n}(z-z_i)$ be monic of degree $n\ge2$, with
$z_i\in\dbar$, the open unit disc.  Show that two of the roots can be joined by a curve of length less than
$2$ lying in the open lemniscate $E=\{z\in\mathbb{C}:|f(z)|<1\}$.
\end{problem}

The open-disc hypothesis is essential for this formulation.  For
$f(z)=z^2-1$, whose roots lie on the unit circle, every continuous path from
$-1$ to $1$ meets the imaginary axis.  There $|f(iy)|=1+y^2\ge1$, so no
such path lies in the open unit lemniscate.  Repeated roots, counted as distinct
occurrences, give a constant path; the substantive case is therefore
squarefree.

\noindent Numbering and current status follow Bloom's Erd\H{o}s problem
catalogue \cite{bloom}, which records the problem as open.
The original source is Problem~5 on printed p.~139 of
Erd\H{o}s--Herzog--Piranian \cite[p.~139]{ehp1958}; the preceding paragraph records the
known input that one component of the lemniscate contains at least two zeros.

Recent work on polynomial lemniscates separates component counts from metric
path questions. Ghosh and Ramachandran characterize the number of components
through critical points and critical values \cite[Lemma~7]{ghosh2023}; for the
binomial family $z^n-a$, the condition $|a|<1$ puts the filled unit lemniscate
in the connected regime. Connectedness alone gives no path-length bound.
Our dynamical terminology is also standard: Sutherland calls
$\dot z=-f(z)/f'(z)$ the continuous Newton flow and observes that $f(z(t))$
moves on a straight radial line \cite[p.~42]{sutherland1992}. We use this
value-space identity, not a claim that the global trajectory graph is a tree.

Two recent manuscripts are relevant.  The 48-page manuscript posted by
\texttt{shtuka} on 24 March 2026
\cite[Theorem~1, p.~1]{march2026} claims the unrestricted
statement.  Its Proposition~12 (p.~16, with proof continuing through p.~30)
supplies the spanning-tree decomposition used in the final proof.  The defect
was located publicly in the problem's discussion thread: on 25 March 2026 Tao
observed that the invocation of Lemma~8 there is unjustified and that the flow
lines need not organise into connected trees, and on 26 March 2026 the
manuscript's author agreed that the statement of Proposition~12 itself, not
only its printed proof, is incorrect, and set the strategy aside.  Section~\ref{sec:gap}
records an independent diagnosis of the same failure through the local
three-ended saddle model, together with possible repairs.  No counterexample
to the proposition is exhibited there.  Pendyala's independent June 2026
preprint \cite[Thm.~1, p.~1]{june2026} proves the degree-four case through a
finite four-point radial lemma and a short polygonal connector.  That is the
degree-four result directly comparable to the root-pair problem. Together
with the cubic theorem established here, it settles these two degrees; it
does not supply the general-degree conclusion. The all-degree estimates below hold in every degree
$n\ge2$, under a hypothesis on the least critical-value modulus or under no
root or threshold hypothesis at all, at the cost of a larger constant or a
weaker containment level.  The quartic theorem does
not close the problem.

We study the Newton flow whose trajectories foliate the lemniscate.

\statusboundary

\begingroup
\footnotesize
\begin{longtable}{@{}
  >{\raggedright\arraybackslash}p{0.34\linewidth}
  >{\raggedright\arraybackslash}p{0.20\linewidth}
  >{\raggedright\arraybackslash}p{0.38\linewidth}@{}}
\toprule
\papertablehead
Statement & Status & Exact boundary \\
\midrule
Erd\H{o}s \#1041 & Open & No proof is claimed. \\
Every monic trinomial $z^n+az^m+b$ & Kernel-checked inequalities; ordinary concatenation & $1\le m<n$ and every zero in the open unit disc; the middle coefficient is unrestricted and $|b|<1$ comes from Vieta.  Lean checks the two radial inequalities and the metric budget; joining the two segments into one path is an ordinary step.  Two intermediate coefficients break the radial mechanism. \\
Separated critical value, $S=2$, $n\ge3$ & Ordinary all-degree theorem & A simple selected saddle with $|1-f(d)/v|\ge2$ at every other critical point; the continuation, injectivity and area arguments are ordinary, and near ties and multiple saddles are excluded.  No existence of a separated critical value is asserted. \\
Low critical value $\mu\le13/25$ & Ordinary all-degree theorem with an exact rational certificate & Squarefree monic, $n\ge2$, least critical-value modulus at most $13/25$; no hypothesis on root locations, component arity or component capacity.  The analytic chain and the certificate are ordinary; nothing here is Lean-checked; the regime $13/25<\mu<1$ is untouched. \\
Unconditional constant factor $(71/10)\mu^{1/n}$ & Ordinary all-degree theorem & Monic, $n\ge2$, no further hypothesis; containment is $\{|f|\le2\mu\}$; the target set is $\{|f|<1\}$, and the constant $71/10$ exceeds the target $2$.  Degree two is the separate exact case with constant $2$. \\
Degree three & Ordinary theorem, complete for that degree & Roots in the open unit disc, listed with multiplicity; the two selected roots are distinct in the squarefree case.  The spoke identity and its norm bound have a Lean companion in the research corpus, outside the pinned formal-source library. \\
Critical-value separation, general $S$ & Ordinary all-degree theorem & A simple nonzero hub and $(2S^2/n)^{1/n}\log(S/(S-1))<1$; the threshold, covering and area arguments remain ordinary; near ties and multiple saddles are excluded. \\
Critical-value budget in every degree & Ordinary theorem; checked analytic assembly & $\sum_j|f(c_j)|^{1/n}\le(n-1)R$ for every $n\ge2$ and $R\ge0$; the repaired aggregate, smoke, and 320 named axiom prints pass. No inverse-ray length estimate is implied. \\
Coefficient energy to path length & Checked & Termwise differentiation and the length integral of the actual power series; the inverse branch and area bound must be supplied. \\
Uniform central free-point region & Checked & Every positive number of points; $|c_i|\le\sqrt{1-e^{-2}}$, with no additional series hypotheses. \\
Newton value equation $w'=-w$ & Checked & Away from critical points, along any trajectory tangent to $-f/f'$. \\
Exponential first integral & Checked & $\tfrac{d}{dt}\bigl(e^{t}f(z(t))\bigr)=0$. \\
Ray separation of critical values & Candidate finite-endpoint completion & Continuity at the endpoints and noncritical Newton evolution on the open interval imply a common positive multiplier; distinct nonzero rays exclude such a connection. No global trajectory producer is supplied. \\
Ray-collision locus & Checked & $\beta=(ra-b)/(1-r)$, $r>0$, $r\ne1$: one real parameter per pair. \\
Quartic case & Cited & Proved in \cite[Thm.~1, p.~1]{june2026}; does not extend to general degree. \\
Translated quartic quotient fibres & Ordinary theorem with a Lean-checked metric kernel & For $f(z)=P((z-h)^q)$, $P$ monic quartic and $q\ge2$; Pendyala supplies the quartic geometry, while Lean checks the root-lift density, exact primitive, and strict endpoint budget. \\
Signed-moment cyclic tetranomials & Lean-checked two-index safe-spoke theorem & For an indexed finite root family of $g(w)=w^m+aw^r+bw^s+c$, an exact signed $L^2$ moment budget selects two distinct indices whose complete spokes are safe; distinct root values require injectivity of the indexing map. \\
Concyclic zeros with $2\rho^n\le1$ & Ordinary proof; finite exact and numerical checks & Distinct-root theorem; not Lean checked; the unrestricted concyclic case remains open. \\
Unrestricted proof of \cite[Theorem~1, p.~1]{march2026} & Proof gap & Proposition~12 uses a false three-ended local saddle block; located publicly by Tao (25 March 2026), conceded by the author at statement level (26 March 2026).  No counterexample is exhibited. \\
Constant-translation ray separation and root retention & Checked & After critical-value injectivity, arbitrary small ray avoidance and an explicit unit-disc margin. \\
Coefficient perturbation and slack stability & Open & Must first create injective critical values and preserve the component, collars and length budget. \\
Reeb decomposition and length fan-in & Open & The two surviving producers. \\
Random search to degree $10$ & Numerical candidate connectors & Grid distances for sampled configurations, with no continuous sublevel certificate on any edge. \\
\bottomrule
\end{longtable}
\endgroup

\section{Trinomials}
\label{sec:trinomials}

Write $E_f=\{z\in\mathbb C:|f(z)|<1\}$.  The first theorem gives a prescribed
path between every pair of zeros of a trinomial, in every degree.

\begin{theorem}[trinomial root connections]
\label{res:trinomial-all-degree}
Let $n,m$ be integers with $1\le m<n$, and let $f(z)=z^n+az^m+b$ have every
zero in $\dbar$.  For every zero $\zeta$, the segment $[0,\zeta]$ lies in
$E_f$.  Consequently any two zeros $\zeta_1,\zeta_2$ are joined in $E_f$ by
the broken line $\zeta_1\to0\to\zeta_2$, of length $|\zeta_1|+|\zeta_2|<2$.
\end{theorem}

The separation assumed here is a condition on critical values, not on the
locations of roots and critical points.  These locations can remain uniformly
separated while two distinct critical values coalesce.  An exact example is
\[
 p(z)=z^4-\frac{z^3}{262144}-2z^2+\frac{3z}{262144}+\frac12,
 \qquad
 p'(z)=4(z+1)(z-3/1048576)(z-1).
\]
Its four roots lie respectively in $(-3/2,-5/4)$, $(-3/4,-1/2)$,
$(1/2,3/4)$, and $(5/4,3/2)$.  Thus the four roots and three critical points
are pairwise more than $1/4$ apart, but the critical values at $-1$ and $1$
differ by only $1/65536<3/65536$.  Root and critical-point separation alone
therefore cannot discharge the value-plane hypothesis of the theorem.

\begin{proof}
Vieta's formula gives $|b|<1$.  At a zero $\zeta$ the root equation
$\zeta^n+a\zeta^m+b=0$ eliminates the middle coefficient:
\begin{equation}\label{eq:trinomial-cancellation}
 f(t\zeta)=b(1-t^m)+\zeta^n(t^n-t^m).
\end{equation}
For $0\le t<1$ the two scalar weights $1-t^m$ and $t^m-t^n$ are nonnegative
and sum to $1-t^n$, so
\[
 |f(t\zeta)|\le|b|(1-t^m)+|\zeta|^n(t^m-t^n)<1-t^n\le1 .
\]
At $t=1$ the value is zero.  Concatenating two such segments gives the length
assertion.
\end{proof}

The coefficient $a$ carries no hypothesis; the root equation removes it before
absolute values are taken.  The conclusion concerns segments to zeros and
makes no assertion that $E_f$ is star-shaped.

\paragraph{The cancellation mechanism.}
\label{subsec:abel-mechanism}
For $f(z)=\sum_{k=0}^nc_kz^k$ and a zero $\zeta$, put
$S_j=\sum_{k=0}^jc_k\zeta^k$.  Finite summation by parts gives
\begin{equation}\label{eq:abel-control-polygon}
 f(t\zeta)=\sum_{j=0}^{n-1}(t^j-t^{j+1})S_j\qquad(0\le t\le1),
\end{equation}
since the coefficient of $c_k\zeta^k$ on the right is $t^k-t^n$ and
$\sum_{k<n}c_k\zeta^k=-c_n\zeta^n$.  The weights
in~\eqref{eq:abel-control-polygon} are nonnegative and sum to $1-t^n$, so the
whole radial segment is safe whenever every partial sum lies in the closed
unit disc.  A trinomial has only two distinct partial sums, $S_j=b$ for $j<m$
and $S_j=-\zeta^n$ for $m\le j<n$, which is exactly the estimate above.

Two intermediate coefficients break this mechanism, and an exact example shows
how.

\begin{example}[an escaping root spoke]\label{ex:sextic-spoke}
Let $0<r<1$ satisfy $r^6>320/327$, and set
\[
 f_r(z)=z^6+\tfrac15r^2z^4-\tfrac15r^4z^2-r^6
       =(z^2-r^2)\bigl(z^4+\tfrac65r^2z^2+r^4\bigr).
\]
Every zero has modulus $r$: for the quartic factor put $z=rw$, so that the
squared roots solve $v^2+(6/5)v+1=0$, whose two roots are complex conjugates
of product one.  Nevertheless
\[
 f_r(r/2)=-\frac{327}{320}\,r^6,
\]
so the segment $[0,r]$ leaves $E_{f_r}$.  This rules out a universal assertion
about every origin-to-zero segment.  It exhibits no counterexample to the
existence of some short connection.
\end{example}

\paragraph{Evidence and exact boundary.}
\label{bdry:trinomial}
The identity~\eqref{eq:abel-control-polygon} is checked as
\href{https://github.com/wcook04/plectis-erdos/blob/0ba585f632fbbbb43af2ee9532d4e83752af3f67/ErdosProblems/Erdos1041/AbelControlPolygon.lean\#L123}{\texttt{abel\_controlPolygon}}, the
constant-term bound as
\href{https://github.com/wcook04/plectis-erdos/blob/0ba585f632fbbbb43af2ee9532d4e83752af3f67/ErdosProblems/Erdos1041/AbelControlPolygon.lean\#L258}{\texttt{norm\_const\_lt\_one\_of\_roots\_lt\_one}},
the radial estimate as
\href{https://github.com/wcook04/plectis-erdos/blob/0ba585f632fbbbb43af2ee9532d4e83752af3f67/ErdosProblems/Erdos1041/AbelControlPolygon.lean\#L219}{\texttt{trinomial\_radial\_norm\_lt\_one}},
and their combination with the metric budget as
\href{https://github.com/wcook04/plectis-erdos/blob/0ba585f632fbbbb43af2ee9532d4e83752af3f67/ErdosProblems/Erdos1041/AbelControlPolygon.lean\#L330}{\texttt{trinomial\_erdos1041\_conclusion}}.
Example~\ref{ex:sextic-spoke} is checked as
\href{https://github.com/wcook04/plectis-erdos/blob/0ba585f632fbbbb43af2ee9532d4e83752af3f67/ErdosProblems/Erdos1041/AbelControlPolygon.lean\#L555}{\texttt{sextic\_guardrail}}.  What the
kernel proves is the pair of radial inequalities and the bound
$|\zeta_1|+|\zeta_2|<2$ on the sum of the two radii.  Assembling those into a
single rectifiable path object, and the passage to the one-dimensional
Hausdorff measure used by the upstream statement of the problem, are ordinary
steps taken here.  Prior art for the trinomial conclusion was searched on
2~September 2026 and none was located; the argument is short enough that it
should be assumed known until a source settles the question.

\section{An unconditional regime: a small least critical value}
\label{sec:low-critical-closure}

The strongest unconditional statement in this note fixes a threshold on one
scalar attached to $f$ and asks nothing else.  Throughout this section
\[
 \mu=\min_{f'(c)=0}|f(c)|
\]
is the least critical-value modulus.

\begin{theorem}[a small least critical value forces a short connector]
\label{res:low-critical-thirteen-twentyfifths}
Let $f$ be squarefree and monic of degree $n\ge2$ with $\mu\le13/25$.  Then
two distinct roots of $f$ are joined inside $\{|f|<1\}$ by a rectifiable
curve of length strictly below $2$.  No hypothesis is placed on the locations
of the roots, on the number of roots in any component, or on the capacity of
any component.
\end{theorem}

\begin{corollary}[scale-free form]
\label{res:low-critical-scale-free}
Every squarefree monic $f$ of degree $n\ge2$ has two distinct roots joined
inside $\{|f|<(25/13)\mu\}$ by a curve of length below
$2\bigl((25/13)\mu\bigr)^{1/n}$.
\end{corollary}

\begin{proof}
Apply Theorem~\ref{res:low-critical-thirteen-twentyfifths} to
$s^{-n}f(sz)$ with $s=\bigl((25/13)\mu\bigr)^{1/n}$, whose least
critical-value modulus is $13/25$, and scale back.  A repeated root occurrence
gives the constant curve, so the squarefree case is the substantive one.
\end{proof}

\paragraph{The failure inequalities.}
Assume no two distinct roots are joined inside $\{|f|<1\}$ by a curve of
length below $2$.  Choose a critical point at level $\mu$ and two descending
inverse arcs into distinct one-root components of $\{|f|<\mu\}$; below $\mu$
each component maps conformally onto the value disc, and distinct local
inverse arcs at the first critical point enter distinct components, since two
inverse images of a nearby regular value in one component would contradict its
degree one.  Their union is a compact connected set containing two roots.
Fix a regular level $t\in(\mu,1)$, let $C_t$ be the ancestor component of that
set, let $k\ge2$ be its root count, and put $x=\log(t/\mu)$ and
$a=\operatorname{Area}(C_t)/\pi$; P\'olya's inequality gives $a\le t^{2/n}<1$.
The component is a Jordan domain, and a Riemann map $\varphi$ turns $f/t$ into
a finite Blaschke product of degree $k$.

For a Riemann map the derivative has Bergman squared norm
$\operatorname{Area}(C_t)$, and the radial restriction of the Bergman kernel
gives $\operatorname{length}(\varphi([0,r]))^2\le a\log(1/(1-r^2))$.  Moving
two preimages to $\pm s$ replaces the kernel integral $-\log(1-r^2)$ by
$4\operatorname{artanh}(s^2)$, so
\[
 \operatorname{length}(\varphi([-s,s]))^2\le4a\operatorname{artanh}(s^2).
\]
A pair of roots at hyperbolic distance $d$ may be placed at $\pm\tanh(d/4)$,
so a connector shorter than $2$ exists as soon as
$\operatorname{artanh}(\tanh^2(d/4))<1/a$.  Failure therefore gives
\begin{equation}\label{eq:lc-separation}
 d(b_i,b_j)\ \ge\ D:=4\operatorname{artanh}\sqrt{\tanh(1/a)},
 \qquad \cosh(D/2)=e^{2/a}.
\end{equation}

Failure also supplies a point $h\in C_t$ whose intrinsic distance to every
root is at least $1$: otherwise the relatively open sets
$\{h:\operatorname{dist}_{C_t}(h,b_j)<1\}$ cover the connected compact set,
and either two of them meet, giving a concatenated path below $2$, or one of
them already contains both roots.  Send $h$ to the origin of the disc, write
$d_j=d(0,b_j)$ and $\lambda(d)=-\log\tanh(d/2)$.  The one-root estimate at $h$
gives $\tanh(d_j/2)\ge\sqrt{1-e^{-1/a}}$, and the Blaschke identity gives the
level budget:
\begin{equation}\label{eq:lc-radius-and-budget}
 \lambda(d_j)\le\frac{\delta(a)}2,\quad
 \delta(a)=-\log\bigl(1-e^{-1/a}\bigr),
 \qquad
 \sum_{j=1}^{k}\lambda(d_j)\ \ge\ x .
\end{equation}

\paragraph{The arity floor.}
The consecutive-gap angular budget that these inequalities suggest is empty:
placing $\lfloor k/2\rfloor$ roots at the minimal radius and the rest at
radius $D$ beyond it, alternating, makes every consecutive law-of-cosines
constraint vacuous while $\sum_j\lambda(d_j)$ grows linearly in $k$.  A root
parked radially behind another costs no angle.  The constraint has to be
imposed at every radius instead.

\begin{lemma}[circle-slice packing]\label{res:circle-slice-packing}
Under~\eqref{eq:lc-separation}, for every $r>0$,
\[
 \sum_{j=1}^{k}w(d_j,r)\le\pi,\qquad
 w(d,r)=\arccos\Bigl(\operatorname{clamp}
   \frac{\cosh d\cosh r-\cosh(D/2)}{\sinh d\sinh r}\Bigr),
\]
where $\operatorname{clamp}$ truncates its argument to $[-1,1]$.
\end{lemma}

\begin{proof}
The open balls $B_j=B_{\mathrm{hyp}}(b_j,D/2)$ are pairwise disjoint, since a
common point would force $d(b_i,b_j)<D$.  Fix $r>0$.  By the hyperbolic law of
cosines the point of the hyperbolic circle of radius $r$ at angle $\theta$
lies in $B_j$ exactly when
$\cosh d_j\cosh r-\sinh d_j\sinh r\cos(\theta-\theta_j)<\cosh(D/2)$, that is
exactly when $|\theta-\theta_j|<w(d_j,r)$ modulo $2\pi$.  So each
$B_j\cap C_r$ is an open arc of angular measure $2w(d_j,r)$, and disjoint arcs
of a circle have total measure at most $2\pi$.
\end{proof}

\begin{theorem}[dual arity floor]\label{res:dual-arity-floor}
Fix radii $r_1,\ldots,r_p>0$ and weights $\sigma_1,\ldots,\sigma_p\ge0$, put
$\Sigma=\sum_i\sigma_i$ and
\[
 U=\sup_{d\ge d_{\mathrm{low}}(a)}
   \Bigl[\lambda(d)-\sum_i\sigma_i\,w(d,r_i)\Bigr],
 \qquad
 \lambda\bigl(d_{\mathrm{low}}(a)\bigr)=\frac{\delta(a)}2 .
\]
If $U>0$, then failure forces $k\ge(x-\pi\Sigma)/U$.
\end{theorem}

\begin{proof}
By~\eqref{eq:lc-radius-and-budget}, Lemma~\ref{res:circle-slice-packing} and
the radius bound,
\[
 x\le\sum_j\lambda(d_j)
  =\sum_j\Bigl[\lambda(d_j)-\sum_i\sigma_i w(d_j,r_i)\Bigr]
    +\sum_i\sigma_i\sum_j w(d_j,r_i)
  \le kU+\pi\Sigma. \qedhere
\]
\end{proof}

Every nonnegative weight vector gives a valid floor, so the weights may be
proposed by any heuristic; the certificate has to bound only $\Sigma$ and $U$.
Both $D(a)$ and $d_{\mathrm{low}}(a)$ decrease in $a$, so a pair
$(\Sigma,U)$ certified at $a'\ge a$ remains valid at $a$, and $k$ is an
integer, so each real floor may be rounded up before the maximum over the
available floors is taken.

\paragraph{From the arity floor to forced area growth.}
At a direction avoiding the finitely many critical-value arguments, lift the
value radius from $0$ to $te^{i\theta}$ from each of the $k$ roots and split
each lift at level $\mu$.  If $A_0$ is the total area of the one-root lobes
below $\mu$ and $\varphi(z)=\sum_l b_lz^l$ maps the disc onto such a lobe,
Cauchy--Schwarz in the radial variable and Parseval give
\[
 \operatorname{mean}_\theta
   \Bigl(\int_0^1|\varphi'(re^{i\theta})|\,dr\Bigr)^2
 \le\sum_{l\ge1}\frac{l^2|b_l|^2}{2l-1}
 \le\sum_{l\ge1}l|b_l|^2
 =\frac{\operatorname{Area}(\text{lobe})}\pi,
\]
so the mean total low lift length is at most $\sqrt{kA_0/\pi}$.  At a regular
intermediate level $u$, the argument principle and the coarea formula give the
perimeter inequality
\begin{equation}\label{eq:lc-perimeter}
 P_C(u)^2\le2\pi k\,u\,A_C'(u),
\end{equation}
because $|dz|=u\,d(\arg f)/|f'|$ on the level curve, its total argument
variation is $2\pi k$, and Cauchy--Schwarz applies.  Integrating
$P_C(u)/(2\pi u)$ from $\mu$ to $t$ bounds the mean high lift length by
$\sqrt{kx(\pi a-A_0)/(2\pi)}$, so some direction has total lift length at most
$M=\sqrt{ka(x+2)/2}$.  Ordering the $k$ boundary endpoints of that direction
cyclically and forming $k$ root connectors, each from two lifts and the
intervening boundary arc, the whole construction lies in
$\overline{C_t}\subset\{|f|<1\}$; under failure every one of them has length
at least $2$, so $2k\le2M+P(t)$.  Applying~\eqref{eq:lc-perimeter} at the
outer level and rearranging gives the area growth inequality
\begin{equation}\label{eq:lc-area-growth}
 a'(x)\ \ge\ \frac1{2\pi^2}
 \Bigl[\,2\sqrt k-\sqrt{2a(x+2)}\,\Bigr]_{+}^{2}
\end{equation}
on every regular interval of the ancestor, with $k$ bounded below by
Theorem~\ref{res:dual-arity-floor}.  At merger levels the ancestor only gains
area, which strengthens the integrated comparison.  P\'olya's inequality caps
$a$ at $1$ until the level reaches $1$, so a trajectory forced past that cap
before $t=1$ contradicts the assumption.

The comparison is integrated cell by cell without assuming that the
independently certified floors form a monotone lookup.  On a cell
$[x_\ell,x_r]$ with trial area $m$ and a certified lower bound $a(x_\ell)\ge
a_\iota$, a value $m$ with $m<a_\iota+(x_r-x_\ell)g(x_\ell,x_r,m)$ is a strict
lower bound for $a(x_r)$: if instead $a(x_r)\le m$, monotonicity of the area
gives $a\le m$ throughout the cell, and integrating~\eqref{eq:lc-area-growth}
with the nonnegative merger jumps contradicts that.

\paragraph{Certificate.}
The threshold is the level at which the forced area passes the cap.  The
companion program
\texttt{experiments/erdos1041\_low\_critical\_path\_certificate.py} evaluates
every accepted inequality in exact rational arithmetic, with directed rounding
on every exponential, logarithm and square root, and with the rational
brackets for $\pi$ and $\log2$ checked from Machin's identity.  The full
replay uses $18$ area-table levels, $14$ radii per dual, $126$ certified
duals, step $1/400$, and the single initial lower area $10^{-6}$.  It
certifies
\[
 X=\frac{635762889599}{1000000000000}<0.6357629,
 \qquad \frac{13}{25}e^{X}<1,
\]
which is Theorem~\ref{res:low-critical-thirteen-twentyfifths}.  The quick mode
certifies $X=664373027131/1000000000000$ and the weaker threshold $51/100$.
The floating optimiser that proposes the dual weights is outside the proof;
only the rational upper bounds on the proposed pairs are used.  Each accepted
pair is a $(\Sigma,U)$ of Theorem~\ref{res:dual-arity-floor}: the quick-mode
pairs at $a=1$ are
$(0.085674,0.045865)$, $(0.126361,0.021829)$, $(0.163157,0.010513)$ and
$(0.208928,0.004318)$, giving $kU+\pi\Sigma$ equal to
$0.4526$, $0.5716$, $0.6808$ and $0.7945$ at $k=4,8,16,32$.  The supremum over
$d$ is taken by worst-first interval refinement, using that $\lambda$ is
decreasing and that $1-w$ is unimodal in $d$ with an interior minimum, so its
minimum on an interval is attained at an endpoint; beyond
$d=\max_ir_i+D/2$ every $w$ vanishes and the bracket is $\lambda$ alone.
The closing inequality is independently checkable: with
$X=635762889599/10^{12}$ and
$e^X\le\sum_{j\le14}X^j/j!+(X^{15}/15!)/(1-X/16)$, rational arithmetic gives
$(13/25)e^X<0.982000386<1$.

\paragraph{Evidence and exact boundary.}
\label{bdry:low-critical-closure}
The analytic chain above is ordinary mathematics.  Its general inputs are the
Riemann mapping theorem, the Bergman kernel, the argument principle, the
coarea formula, and P\'olya's area inequality
$\operatorname{Area}\{|f|\le t\}\le\pi t^{2/n}$
\cite[printed pp.~280--282]{polya1928}.  The exact rational certificate is a
kernel-free replayable computation, and no part of this theorem is
Lean-checked or independently reviewed.  Prior art for the assembled statement
is unassessed; the hyperbolic slice inequality follows from disjointness of
the balls and the law of cosines and should be assumed known.  The public
source of record is
\texttt{research\_corpus/Erdos1041/LowCriticalPathCertificate.md}, which
carries the full analytic chain, the certified dual table and the replay
route.  The exact remaining obligation is that this threshold argument does
not cover $13/25<\mu<1$; other results in this note intersect that range, so
the globally open set is not obtained by subtracting one scalar regime.
Fixed-degree arguments remain stronger at $n=4$ and $n=5$, where the
corresponding thresholds are $61/100$ and $139/250$; from $n=6$ on the
all-degree constant $13/25$ is the better statement.  Erd\H{o}s~\#1041 remains
open.

\section{An unconditional all-degree constant-factor path theorem}
\label{sec:constant-factor}

Dropping the threshold entirely costs a constant.  For a monic degree-$n$
polynomial put
\[
 K_t=\{z:|f(z)|\le t\},\qquad
 \mu=\min_{f'(c)=0}|f(c)|,\qquad \rho=\mu^{1/n}.
\]

\begin{theorem}[unconditional constant-factor path]
\label{res:constant-factor-path}
For every monic polynomial $f$ of degree $n\ge2$, two zero occurrences are
joined by a possibly degenerate path of length at most
\[
 \frac{71}{10}\,\rho
\]
inside $K_{2\mu}$.  If $f$ is squarefree, their locations are distinct.
If every root lies in the open unit disc and $\mu\le1/2$, the construction
may be chosen inside $\{|f|<1\}$ with length at most $5.7$.
\end{theorem}

The constant is the rationally certified specialization of a two-parameter
bound.  If the selected working component contains $k\ge2$ roots, then for
every $r\in(0,1)$ and $\lambda>1$ the proof constructs a path in
$K_{\lambda\mu}$ whose length is at most
\[
 \sqrt{\frac2k}\left(
   \frac{\sqrt2\,r}{(1-r)^2}
   +\lambda^{1/n}\left(
      \sqrt{\log(\lambda/r)}+
      \frac{\pi}{\sqrt{\log\lambda}}
    \right)
 \right)\rho.                                      \tag{CF}
\]
For $n\ge3$, the choice $\lambda=2$, $r=3/20$ makes the bracket less
than $71/10$: the largest case is $n=3$, where exact rational bounds give
$66517563/9392500<71/10$.  Degree two has the exact root-segment bound
$2\rho$.

\begin{proof}[Proof architecture]
Choose a critical point at level $\mu$ and a component at a regular level
$t\in[\mu,\lambda\mu]$ that contains the corresponding first merge.  If a
regular component $C'$ at level $\sigma$ contains $k'$ roots, winding and
Cauchy--Schwarz give the local perimeter estimate
\[
 \mathcal H^1(\partial C')^2
 \le2\pi k'\,\sigma\,\frac{d}{d\sigma}
      \operatorname{Area}(K_\sigma\cap C').         \tag{CF1}
\]
P\'olya's area--capacity inequality
$\operatorname{Area}(K_T)\le\pi T^{2/n}$ then selects $t$ with a
degree-free boundary-length bound.

For an argument avoiding the finitely many critical-value rays, lift the
radial value segment from each root to $\partial C_t$.  Split each lift at
level $r\mu$.  Koebe distortion controls the low pieces.  The exact
precritical conformal-radius energy inequality
\[
 \sum_i\frac{\mu^2}{|f'(z_i)|^2}\le\rho^2          \tag{CF2}
\]
aggregates those pieces without paying once per root.  Coarea and~(CF1)
control the high pieces on average over the argument.  Choose one argument
realising that average, order the lift endpoints cyclically on
$\partial C_t$, and use the shortest adjacent boundary arc.  Averaging the
two lifts and the boundary arcs over the $k$ adjacent pairs gives~(CF).
The strict unit-sublevel statement chooses the regular level below the top of
the window when $\lambda\mu\le1$.
\end{proof}

The bracket~(CF) retains two quantities that the constant $71/10$ discards,
namely the root count $k$ of the working component through the factor
$\sqrt{2/k}$, and the capacity of that component through the area input.
Keeping either one turns the constant-factor theorem into the target
conclusion on an explicit region.

\begin{corollary}[high-arity closure]
\label{res:constant-factor-arity}
Let $f$ be monic with every root in the open unit disc, let $c_*$ be a
critical point with $|f(c_*)|=\mu$, and let $k_0$ be the number of roots,
counted with multiplicity, in the component of $K_\mu$ containing $c_*$.
Then Erd\H{o}s~\#1041 holds for $f$ in each of the three cases
\[
 \mu\le\tfrac12\ \text{and}\ k_0\ge17,\qquad
 \mu\le\tfrac14\ \text{and}\ k_0\ge12,\qquad
 \mu\le\tfrac18\ \text{and}\ k_0\ge10 .
\]
\end{corollary}

\begin{proof}
Every working component $C_t$ in the proof of
Theorem~\ref{res:constant-factor-path} contains the first-merge component, so
$k\ge k_0$.  For the first case take $\lambda=2$, $r=13/100$; then
$(2\mu)^{1/n}\le1$ and $\rho\le1$, and the exact bounds
$\sqrt2<283/200$, $\sqrt{\log(200/13)}<5/3$ and
$\pi/\sqrt{\log2}<(22/7)/(104/125)=1375/364$ make the bracket in~(CF) at most
$15668813/2755116$, whose square is $34-12570881068535/7590664173456<34$.
Hence $k_0\ge17$ gives squared length below $(2/17)\cdot34=4$.  For the second
case take $\lambda=4$, $r=3/25$: the bracket is below $6075221/1273888$, whose
square is $24-2038665078215/1622790636544<24$, and $2/k_0\le1/6$ gives squared
length below $4$.  For the third take $\lambda=8$, $r=11/100$: the bracket is
below $55629121/12475575$, whose square is
$20-18200328379859/155639971580625<20$, and $2/k_0\le1/5$ again gives squared
length below $4$.  In each case $\lambda\mu\le1$, so Lemma-level freedom in
the choice of the regular level $t$ keeps $t<1$ and the containment strict.
\end{proof}

\begin{corollary}[capacity-defect closure]
\label{res:constant-factor-capacity}
Keep the hypotheses of Corollary~\ref{res:constant-factor-arity} with
$\mu\le1/2$, let $C$ be the component of $\{|f|<2\mu\}$ containing $c_*$, and
put $\kappa=\operatorname{cap}(\overline C)/(2\mu)^{1/n}$.  If
$\kappa\le\tau_{k_0}$, where
\[
 \tau_k=\frac{\sqrt{2k}-A}{B},\qquad
 A=\frac{283}{3610},\qquad B=\frac{52029}{9100},
\]
then Erd\H{o}s~\#1041 holds for $f$.  In particular $\kappa\le1/3$ suffices at
every arity, and the rational cutoffs
$2/5,\,12/25,\,1/2,\,7/12,\,16/25,\,2/3,\,7/10$ suffice at
$k_0=3,\ldots,9$, rising to $39/40$ at $k_0=16$.
\end{corollary}

\begin{proof}
Replace the global input $\operatorname{Area}(K_T)\le\pi T^{2/n}$ in the proof
of Theorem~\ref{res:constant-factor-path} by the exact component form
$\operatorname{Area}(C)\le\pi\operatorname{cap}(\overline C)^2
 =\pi\kappa^2(2\mu)^{2/n}$.  Only the high-lift and boundary terms acquire the
factor $\kappa$; the low-lift term~(CF2) is unchanged.  Taking $\lambda=2$,
$r=1/20$, and the exact bounds above together with
$\log40<12641/3402<(97/50)^2$, gives
$\operatorname{length}<\sqrt{2/k_0}\,(A+B\kappa)$.  The definition of
$\tau_{k_0}$ makes the right side at most $2$, and for each displayed rational
$q_k$ integer arithmetic gives $(A+Bq_k)^2<2k$.  Discarding
$\sqrt{2/k_0}\le1$ altogether gives the uniform cutoff $\kappa\le1/3$.
\end{proof}

By the component-capacity formula, $\kappa=e^{-\Sigma/n}$ with $\Sigma$ the sum
of exterior Green function values at the roots excluded from $C$.  A remaining
counterexample with $\mu\le1/2$ therefore has $k_0\le16$ and exterior Green
defect $\Sigma<n\log(40/39)$ at the top arity, so it must be simultaneously
low-arity and nearly round.

The level window $[\mu,\lambda\mu]$ is the only loss in
Theorem~\ref{res:constant-factor-path} that the mechanism itself demands, and
removing it isolates one clean extremal problem.  At the critical level
itself~(CF1) fails, because the area derivative diverges logarithmically at a
critical level.

\begin{problem}[one-lobe perimeter]\label{prob:conjecture-p}
Let $C$ be a connected component of $\{|f|\le\sigma\}$ containing exactly one
root of $f$.  Prove or refute
\[
 \mathcal H^1(\partial C)\le\sqrt2\,\varpi\,\sigma^{1/n}
 =\sqrt2\,\varpi\operatorname{cap}\{|f|\le\sigma\},
 \qquad
 \sqrt2\,\varpi=\frac{\Gamma(1/4)^2}{2\sqrt\pi}=3.70814935\ldots,
\]
where $\varpi$ is the lemniscate constant.
\end{problem}

The constant is attained: for $f=z^2-d^2$ one has $\mu=d^2$, $\rho=d$, and the
one-root lobe boundary $|z^2-d^2|=d^2$ is one loop of Bernoulli's lemniscate,
of length $d\,\Gamma(1/4)^2/(2\sqrt\pi)$ by the substitution
$\varphi(w)=d\sqrt{1+w}$ and $\Gamma(1/4)\Gamma(3/4)=\pi\sqrt2$.  On
$z^n-r^n$ the corresponding closed form
$(2^{1/n}/n)\sqrt\pi\,\Gamma(1/(2n))/\Gamma(1/(2n)+1/2)$ decreases from
$3.7081$ at $n=2$ to $2$ as $n\to\infty$.  Problem~\ref{prob:conjecture-p} is
the one-root per-component analogue of the Erd\H{o}s--Herzog--Piranian
lemniscate-length problem, with a different extremal: their per-root constant
is $2$, attained asymptotically by $z^n-1$, while the one-root maximum here is
attained by a quadratic.

\begin{theorem}[conditional removal of the level window]
\label{res:conjecture-p-consumer}
If the bound of Problem~\ref{prob:conjecture-p} holds with constant $\beta$
for every $\sigma\le\mu$, then two roots of $f$ are joined inside $K_\mu$
by a path of length at most
\[
 \Bigl(\frac{8x}{(1-x)^2}+\frac\beta\pi\log\frac1x+\beta\Bigr)\rho
 \qquad(0<x<1),
\]
which is about $7.4\rho$ at $\beta=\sqrt2\varpi$ and $x=1/10$.
\end{theorem}

\begin{proof}
At level $\mu$ exactly, $c_*$ lies on the boundary of two one-root lobes
$U_a,U_b$, so it suffices to bound the intrinsic distance from $a$ to $c_*$
inside $\overline{U_a}$.  Lift from $a$ along a ray to $\partial U_a$ and
travel along that Jordan curve to $c_*$, so that
$d(a,c_*)\le\ell_a(\theta)+\tfrac12\mathcal H^1(\partial U_a)$.  The mean lift
is $(2\pi)^{-1}\int_0^\mu P_a(\sigma)\,d\sigma/\sigma$ with
$P_a(\sigma)=\mathcal H^1(\partial U_a(\sigma))$.  Bound $P_a$ by the
distortion estimate for $\sigma\le x\mu$ and by the hypothesis for
$\sigma\ge x\mu$, using $\mathrm{cr}_a\le4\rho$ and
$P_a(\sigma)\le\beta\sigma^{1/n}\le\beta\rho$.  Doubling for the two lobes and
adding $\tfrac12\beta\rho$ twice gives the stated bound.
\end{proof}

\paragraph{Evidence and exact boundary.}
\label{bdry:constant-factor}
Theorem~\ref{res:constant-factor-path} is ordinary mathematics with no
Lean-checked part.  It uses the capacity identity, P\'olya's area inequality,
Cauchy--Schwarz, the univalent-map distortion estimate, and a finite graph
exit lemma; no external paper theorem is imported into it.  The public source
of record is
\texttt{research\_corpus/Erdos1041/UnconditionalConstantFactorBound.md}, which
carries every rational verification quoted above.  Its
containment level is $2\mu$ and its constant is $71/10$, so the theorem itself
reaches neither half of the target, and the sharp constant $2$ remains open.
Corollaries~\ref{res:constant-factor-arity} and~\ref{res:constant-factor-capacity}
do reach the target on their stated regions; each rational cutoff is an integer
inequality of the form $(A+Bq)^2<2k$, checked exactly.  The complementary
cells, low first-merge arity together with high component capacity, stay open.
Problem~\ref{prob:conjecture-p} is a conjecture and is measured only: $120$
configurations over degrees $2$ through $10$ and five families give a maximum
of $\mathcal H^1(\partial U)/(\sqrt2\varpi\rho)$ equal to $0.9999997747$,
attained at quadratics, with the constructed path at $5.07\rho$ against the
proved $7.1\rho$; no adversarial search has been run against it, so the
reported maximum is a basin record.  The
adjacent classical literature on lemniscate length measures a different
object, namely the arclength of the level curve $\{|p|=1\}$, which is
Erd\H{o}s~\#114: Borwein's $8\pi en$
\cite{borwein1995}, Eremenko and Hayman's $9.173n$ \cite{eremenkohayman1999},
Fryntov and Nazarov's asymptotically sharp $2n+o(n)$ \cite{fryntovnazarov2008},
and Tao's resolution for large $n$ \cite{tao2025}.  A bound $74n^2$ is
attributed to Pommerenke in the same family, and this note records no located
reference for it.  None of these bounds a
root-to-root path, and a dated search on 2 September 2026 located no
constant-factor root-pair bound in arbitrary degree.

\section{Degree three}
\label{sec:degree-three}

\begin{theorem}[the cubic case]
\label{res:degree-three}
Let $f(z)=\prod_{j=1}^{3}(z-z_j)$ with $|z_j|<1$, the roots listed with
multiplicity.  Then two listed root occurrences are joined inside
$\{|f|<1\}$ by a polygonal path of length strictly below $2$.  If $f$ is
squarefree the two are distinct.
\end{theorem}

\noindent\href{https://github.com/wcook04/plectis-erdos/blob/bb4e24651b37cd096d3749ef299c3317930b3a3b/ErdosProblems/Erdos1041/PaperCubicCompletion.lean#L297}{The complete cubic construction is checked in Lean, including its polygonal path, strict variation bound, and repeated-root case.}


\begin{proof}
Multiple roots give a constant path, so assume $f$ squarefree.  A component
containing two roots contains a critical point, so the
Erd\H{o}s--Herzog--Piranian component lemma \cite[p.~139]{ehp1958} supplies a
critical point with critical-value modulus below one.

Suppose first that $f'$ has two distinct zeros.  Choose a critical point $c$
minimising $|f(c)|$, write the other one as $c+\delta$, and put $v=f(c)$, so
$0<|v|<1$.  Monicity gives the exact expansion
$f(c+d)=d^3-\tfrac32\delta d^2+v$.  Choose $\alpha$ with $\alpha^3=v$ and set
$b=\delta/\alpha$; dividing by $v$ gives
$f(c+\alpha w)/v=P_b(w)$ with
\[
 P_b(w)=w^3-\tfrac32bw^2+1 .
\]
At the other critical point $f(c+\delta)=v(1-b^3/2)$, so minimality of $|v|$
is exactly the hypothesis $|1-b^3/2|\ge1$.

Under that hypothesis $P_b$ has at least two zeros, with multiplicity, in the
closed unit disc.  In the strict region a unit-circle zero $w$ would give
$b=\tfrac23(w+w^{-2})$, and with $w^3=e^{2iu}$ and $c_u=\cos u$ a direct
calculation gives $|1-b^3/2|^2=1+\tfrac{64}{729}c_u^4(16c_u^2-27)\le1$, an
equality only at $c_u=0$, where $b=0$.  So the strict region has no
unit-circle zero.  The radial deformation $b\mapsto sb$, $s\ge1$, stays in
the strict region, since $z=b^3/2$ and $|1-z|^2>1$ give
$|1-s^3z|^2-1=s^3(s^3|z|^2-2\Re z)>0$.  For large $s$, Rouch\'e's theorem on
$|w|=1$ compares $P_{sb}$ with $-\tfrac32sbw^2$, whose modulus
$\tfrac32s|b|$ exceeds $2\ge|w^3+1|$, so $P_{sb}$ has exactly two zeros in
the open unit disc; no zero crosses the circle along the deformation.  The
equality case follows by letting $s$ decrease to one and using continuity of
the root multiset.

If $P_b(w)=0$ and $0\le t\le1$, then
$P_b(tw)=1-t^2-t^2(1-t)w^3$, so $|w|\le1$ gives
$|P_b(tw)|\le(1-t^2)+t^2(1-t)|w|^3\le1-t^3\le1$.  The whole segment from $0$
to $w$ is therefore safe.  Selecting two normalized roots $w_1,w_2$ with
$|w_i|\le1$, the two segments from $c$ to $c+\alpha w_i$ lie in
$\{|f|\le|v|\}\subset\{|f|<1\}$ and have combined length at most
$2|\alpha|=2|v|^{1/3}<2$.

If instead $f'$ has a double zero $c$, then $f(c+d)=d^3+v$ and the roots are
$c+\alpha$, $c+\alpha\omega$, $c+\alpha\omega^2$.  Averaging their squared
moduli gives $|c|^2+|\alpha|^2<1$, so $|\alpha|<1$; any two radial spokes
have total length $2|\alpha|<2$, and along either spoke the value has modulus
$|v|(1-t^3)<1$ away from the root endpoint.
\end{proof}

\paragraph{Evidence and exact boundary.}
\label{bdry:degree-three}
The cubic reduction, zero-count lemma and path assembly are now checked in
\texttt{ErdosProblems/Erdos1041/PaperCubicCompletion.lean}, with endpoint
\texttt{cubic\_paper\_complete}; the supporting root count is in
\texttt{PaperCubicSchur.lean}. The
\href{https://github.com/wcook04/plectis-erdos/blob/3be82b1a7340284aea72e9a5c8493cb020843921/ErdosProblems/Erdos1041/PaperCubicMonic.lean#L32}{monic-cubic endpoint} also discharges the root enumeration:
monicity and degree three supply the factorisation over $\mathbb C$, with
multiplicities retained.  The connector therefore starts from the polynomial
and its open-disc root condition alone; squarefreeness gives distinct root
endpoints. The earlier research-corpus presentation is
historical working material, rather than the present proof boundary.  Pendyala proves the
degree-four case \cite[Thm.~1, p.~1]{june2026}; whether degree three was
already recorded on the Erd\H{o}s Problem~\#1041 discussion thread or in
\cite{ehp1958} has not been adjudicated, so the priority question here is
open.  The argument uses the sparse normalized form $w^3-\tfrac32bw^2+1$,
and a higher-degree proof needs a replacement for that sparsity.

\section{An all-degree critical-value separation theorem}
\label{sec:critical-value-separation}

The first result is the strongest direct parent-theorem regime in this note.
It is not a statement about a sampled family or a fixed degree: the condition
is an exact inequality on the other critical values after normalising one
simple first-merge hub.

\begin{theorem}[critical-value separation]
\label{res:critical-value-separation}
Let $P$ be a polynomial of degree $n\ge3$ whose leading coefficient has
modulus one.  Suppose
\[
 P(0)=1,\qquad P'(0)=0,\qquad P''(0)\ne0,
\]
and, for some $S>1$, every other critical point $d\ne0$ satisfies
\[
 |1-P(d)|\ge S.
\]
If $Z$ is either local solution at the saddle of
\[
 P(Z(\xi))=1-\xi^2,\qquad Z(0)=0,
\]
then $Z$ continues holomorphically and injectively to
$|\xi|<\sqrt S$.  Its endpoints $Z(-1)$ and $Z(1)$ are distinct roots, the
curve $Z([-1,1])$ lies in $\{|P|\le1\}$, and their resolved inverse-ray
connector satisfies
\[
 \int_{-1}^{1}|Z'(\xi)|\,d\xi
 \le 2\Bigl(\frac2n\Bigr)^{1/(2n)}S^{1/n}
     \sqrt{\log\!\frac{S}{S-1}}.       \tag{4}
\]
In particular the connector has length strictly below $2$ whenever
\[
 \Bigl(\frac{2S^{2}}{n}\Bigr)^{1/n}\log\!\frac{S}{S-1}<1. \tag{5}
\]
\end{theorem}

\begin{proof}
The simple saddle has local form $P(z)=1+Az^2+O(z^3)$ with $A\ne0$, so the
substitution $P(Z(\xi))=1-\xi^2$ resolves it into two local holomorphic
sheets.  Normalize the algebraic curve
$X=\{(\xi,z):P(z)=1-\xi^2\}$.  A finite ramification point of its projection
to the $\xi$-plane can occur only at a critical point $d$ of $P$, where
$\xi^2=1-P(d)$.  The node over $(0,0)$ is replaced in the normalization by
the two resolved local points, and the separation hypothesis puts every
other ramification point on or outside $|\xi|=\sqrt S$.  Polynomial
properness makes the normalized projection finite and proper.  Its inverse
image over $D_{\sqrt S}$ is therefore an unramified finite cover.  Since the
disc is simply connected, every connected cover component maps
biholomorphically to it; in particular either resolved local sheet continues
as a single holomorphic branch $Z$ throughout the disc.

This branch is also injective as a map into the $z$-plane.  Indeed,
$Z(\xi_1)=Z(\xi_2)$ implies $\xi_1^2=\xi_2^2$.  The involution
$\iota(\xi,z)=(-\xi,z)$ exchanges the two distinct normalized points over
$(0,0)$ and hence exchanges their two cover components.  If
$\xi_2=-\xi_1\ne0$ and $Z(\xi_1)=Z(\xi_2)$, one component would meet its
$\iota$-image; connected cover components are disjoint, so they would have to
coincide, contradicting the two distinct points over $0$.  Thus
$\xi_1=\xi_2$.

Write $Z(\xi)=\sum_{k\ge1}a_k\xi^k$.  For
$1<R<\sqrt S$, injectivity and the area formula give
\[
 \pi\sum_{k\ge1}k|a_k|^2R^{2k}
   =\operatorname{Area}(Z(D_R)).                    \tag{6}
\]
The image is bounded and open, and $Z$ is holomorphic on a neighbourhood of
$\overline{D_R}$, so the area of $U_R:=Z(D_R)$ may be computed by the change
of variables $z=Z(\xi)$, whose real Jacobian is $|Z'|^2$.  Differentiating
$P(Z(\xi))=1-\xi^2$ gives $P'(Z(\xi))Z'(\xi)=-2\xi$, and therefore the exact
Dirichlet integral
\[
 \int_{U_R}|P'(z)|^2\,dA(z)
 =\int_{D_R}\bigl|P'(Z(\xi))Z'(\xi)\bigr|^2\,dA(\xi)
 =4\int_{D_R}|\xi|^2\,dA(\xi)=2\pi R^4.            \tag{7}
\]
Write $P=e^{i\alpha}\widetilde P$ with $\widetilde P$ monic; this changes
neither $|P'|$ nor any area.  Crane's sharp lower bound for the area of a
polynomial image counted with multiplicity \cite[Theorem~3]{crane} states that
$\int_K|\widetilde P'|^2\,dA\ge n\pi(\operatorname{Area}(K)/\pi)^n$ for every
measurable $K$.  Applied to $K=U_R$ and combined with~(7), it gives
\[
 \operatorname{Area}(U_R)
 \le \pi\Bigl(\frac{2R^4}{n}\Bigr)^{1/n}
 = \pi\Bigl(\frac2n\Bigr)^{1/n}R^{4/n}.            \tag{8}
\]
Cauchy--Schwarz applied to (6) and (8) now gives
\[
 \sum_{k\ge1}|a_k|
 \le \Bigl(\frac2n\Bigr)^{1/(2n)}R^{2/n}
      \sqrt{\sum_{k\ge1}\frac1{kR^{2k}}}
 = \Bigl(\frac2n\Bigr)^{1/(2n)}R^{2/n}
   \sqrt{\log\!\frac{R^2}{R^2-1}}.                  \tag{9}
\]
The \href{https://github.com/wcook04/plectis-erdos/blob/0d34630e1cc9d2b1ac6edfa7cfdba83b31bdc8fe/ErdosProblems/Erdos1041/PowerSeriesDerivative.lean#L120}{checked coefficient-to-length theorem} formalises the passage from the subradius energy bounds to the actual power-series derivative. Termwise integration on $[-1,1]$ gives
$\int_{-1}^{1}|Z'|\le2\sum_{k\ge1}|a_k|$.  Letting
$R\nearrow\sqrt S$ proves (4).  Injectivity makes $Z(-1)$ and $Z(1)$
distinct, while $P(Z(\xi))=1-\xi^2\in[0,1]$ on the real segment proves the
claimed containment.  Squaring (4) gives (5).
\end{proof}

\begin{corollary}[separation two suffices in every degree]
\label{res:critical-value-thresholds}
Condition~\emph{(5)} holds at $S=2$ for every integer $n\ge3$.
Consequently, if $f$ is monic with roots in the open unit disc, $c$ is a
simple critical point with $v=f(c)\ne0$, and every other critical point
satisfies $|1-f(d)/v|\ge2$, then two roots of $f$ are joined inside
$\{|f|<1\}$ by a curve of length strictly below $2$.
\end{corollary}

\begin{proof}
At $S=2$ the left side of~\emph{(5)} is $B_n=(8/n)^{1/n}\log2$.  For
$3\le n\le8$ the base $8/n$ is at least one and decreases with $n$, and the
exponent $1/n$ decreases as well, so $(8/n)^{1/n}\le(8/3)^{1/3}$; for $n\ge8$
the base is at most one, so $(8/n)^{1/n}\le1\le(8/3)^{1/3}$.  Next,
\[
 e^{7/10}>1+\frac7{10}+\frac{49}{200}+\frac{343}{6000}
        =\frac{12013}{6000}>2,
\]
because every term of the exponential series is positive, so $\log2<7/10$.
Cubing the resulting bound on $B_n$ gives
\[
 \Bigl[\Bigl(\frac83\Bigr)^{1/3}\frac7{10}\Bigr]^{3}
 =\frac83\cdot\frac{343}{1000}=\frac{343}{375}<1,
\]
so $B_n<1$ in every degree $n\ge3$.

For the polynomial $f$, set $z=c+\rho e^{i\theta}w$ with
$\rho=|v|^{1/n}$ and choose $\theta$ so that $P(w)=f(z)/v$ has
unit-modulus leading coefficient.  The theorem gives a connector of length
strictly below $2|v|^{1/n}$ inside $\{|f|\le|v|\}$.  At $S=2$, separation
gives $|f(d)|\ge(S-1)|v|=|v|$ at every other critical point, so $|v|=\mu$ is
the least critical-value modulus.  If the roots lie in a disc of radius
$R<1$, the resultant identity and Fekete's Vandermonde bound give, with
critical points counted with multiplicity,
\[
 \mu^{n-1}
 \le \prod_{f'(d)=0}|f(d)|
 =\frac{|\operatorname{Disc}(f)|}{n^n}
 =\frac{\prod_{i<j}|z_i-z_j|^2}{n^n}
 \le R^{n(n-1)}.
\]
The last step writes $z_i=z_0+Ru_i$ with $|u_i|\le1$ and applies Hadamard's
inequality to the Vandermonde matrix $(u_i^{j-1})$, each of whose rows has
Euclidean norm at most $\sqrt n$.  Thus $|v|=\mu\le R^n<1$, so both the length
and containment are strict at the target scale.
\end{proof}

The threshold $S=2$ is sufficient and no optimality is asserted.  As $n$ grows
the first factor of \emph{(5)} tends to one, so the limiting cutoff is the
solution of $\log(S/(S-1))=1$, namely $S=e/(e-1)=1.5820\ldots$; at $n=3$ the
numerical cutoff is near $1.903$.  For orientation, the normalised length bound at
$S=2$ is about $1.9608$ in degree three, $1.8158$ in degree four, and
$1.7452$ in degree five.  The polynomial $P(z)=1-3z^2+z^3$ satisfies the
hypotheses with $S=4$: its other critical point is $2$, and $1-P(2)=4$.

\paragraph{Evidence and exact boundary.}
\label{bdry:critical-value-separation}
The analytic continuation, monodromy, area formula, change of variables and
external area inequalities in Theorem~\ref{res:critical-value-separation} are
ordinary mathematics.  Two external theorems enter, both used as stated.
P\'olya's area inequality \cite[printed pp.~280--282]{polya1928} bounds the
area of a polynomial sublevel set; Crane's Theorem~3 \cite{crane} bounds from
below the area of a polynomial image counted with multiplicity, sharply, with
equality exactly when the set is a disc and the polynomial has a single
critical value at its centre.  The route through P\'olya alone, applied to
$P-1$ on $\{|P-1|<R^2\}$, already gives
$\operatorname{Area}(U_R)\le\pi R^{4/n}$ and the length bound
$2S^{1/n}\sqrt{\log(S/(S-1))}$; the factor $(2/n)^{1/(2n)}$ recorded in~(4) is
the additional gain from the exact Dirichlet integral~(7), and it is what
carries $S=2$ down to degree three.  The coefficient-energy estimate, termwise
differentiation and length integral are Lean-checked for the actual power
series; the companion symbolic replay checks the branch-value algebra and the
boundary example $P(z)=1-3z^2+z^3$.  Constructing the univalent branch remains
an ordinary input to the full theorem.

The method stops when a
second critical value enters the resolved disc, and it assumes the selected
saddle is simple.  It therefore leaves the near-tie and multiple-saddle
strata, makes no sharpness claim for the threshold $S=2$, asserts no existence
of a separated critical value, and does not prove an unrestricted
admissible-hub selector, a COVER theorem, or Erd\H{o}s~\#1041.

The literature disposition for Dubinin's Theorem~1 is settled here.  The
source is \cite{dubinin}, Theorem~1 on printed page~85 of the POMI original.
It concerns a holomorphic function that gives a full $n$-fold covering of an
annulus $t_1<|w|<t_2$, and, in the notation of that theorem, with $E$ its
explicitly defined complementary set, it states
\[
 \Bigl(\frac{t_2}{t_1}\Bigr)^{2/n}\le\frac{m(E\cup D)}{m(E)}.
\]
The full covering and the explicit complementary set are hypotheses of the
theorem.  Tao cited it on the Erd\H{o}s Problem~\#1041 discussion page on
25 March 2026 for the relative scaling factor $s^{2/n}$ between the areas of
two nested sublevel sets, which is the reading of the displayed inequality in
which $E$ and $E\cup D$ are those two sublevel sets and the covering
hypothesis holds.  No step of this note uses that relative inequality.  The
area estimate~(8) above uses Crane's absolute image inequality together with
the exact Dirichlet integral~(7), and the absolute sublevel inequality
$\operatorname{Area}\{|P|<T\}\le\pi T^{2/n}$ for a polynomial with
unit-modulus leading coefficient, taken from P\'olya
\cite[printed pp.~280--282]{polya1928}, is the area input to
Theorem~\ref{res:low-critical-thirteen-twentyfifths} and to
Theorem~\ref{res:constant-factor-path}.  Dubinin's theorem is a
neighbouring result on lemniscate areas under a covering hypothesis.  It
bounds areas and supplies no curve joining two roots, so it leaves every
statement of this note unchanged.  Corollary~\ref{res:critical-value-thresholds}
states $S=2$ in every degree $n\ge3$; earlier forms of this threshold recorded
$n\ge7$ and then $n\ge6$, and the exact image-area step~(7)--(8) supersedes
both.

\section{Three exact solved polynomial families}
\label{sec:solved-polynomial-families}

The next three theorems are unconditional subcases of the path problem, not
proxies for the unrestricted conjecture.  They are ordered by mathematical
signal: first an all-degree sharp theorem with equality configurations, then a
complete degree-five sparse family, and finally an infinite tower of translated
quotient fibres.  The date of their formalization has no bearing on that order.
In each case a displayed finite proposition is the exact Comparator-selected
Lean endpoint.  The ensuing normalization and path theorem is ordinary
mathematics and is stated separately.

\subsection{Sharp collinear zeros}
\label{subsec:sharp-collinear-chebyshev}

Put
\[
 r_n=\cos\frac{\pi}{2n},
 \qquad C_n=\frac{1}{2^{n-1}r_n^n}.
\]

\begin{theorem}[checked Chebyshev endpoint]
\label{prop:sharp-collinear-chebyshev-comparator}
Let $m\ge0$, let $p\in\mathbb R[X]$ be monic of degree $m+2$, and let
\[
 -1<c_0<\cdots<c_m<1,\qquad |c_i|\le1.
\]
Suppose $p(-1)=p(1)=0$ and
$p(c_i)p(c_{i+1})<0$ for $0\le i<m$.  Then
\[
 \min_{0\le i\le m}|p(c_i)|\le C_{m+2}.
\]
\end{theorem}

This is exactly the statement selected as
\href{\repobase/ExternalVerification1041SolvedFamilies/Solution.lean\#L74}{the formal Chebyshev endpoint}.
It is the kernel-facing alternation endpoint, not yet the geometric theorem.

\begin{theorem}[sharp collinear diameter theorem]
\label{thm:sharp-collinear-diameter}
Let $f$ be a monic polynomial of degree $n\ge2$ whose zero occurrences are
collinear, and let $D$ be their diameter.  Some two adjacent zero occurrences
are joined by a segment of length at most $D$ on which
\[
 |f(z)|\le
 \frac{(D/2)^n}{2^{n-1}\cos^n(\pi/(2n))}.          \tag{9}
\]
The constant in \emph{(9)} is best possible in every degree.  Equality is
attained by affine images of the zeros of $T_n$ whose extreme zeros have
distance $D$.
\end{theorem}

\begin{proof}
A repeated zero gives the constant path, so assume the zero values are
distinct.  A rigid motion sends their line to the real axis, their midpoint
to the origin, and their extremes to $\pm D/2$.  With $R=D/2$, the normalized
polynomial
\[
 q(w)=R^{-n}e^{-in\theta}f(m+Re^{i\theta}w)
\]
is monic with real zeros
$-1=y_1<\cdots<y_n=1$, and
$|f(m+Re^{i\theta}w)|=R^n|q(w)|$.

Compare $q$ with the monic endpoint-normalized Chebyshev polynomial
\[
 q_*(x)=\frac{T_n(r_nx)}{2^{n-1}r_n^n}.
\]
Both $q$ and $q_*$ vanish at $\pm1$, and $|q_*|\le C_n$ on
$[-1,1]$.  For each gap $[y_i,y_{i+1}]$, choose $c_i$ at which
$|q|$ is maximal.  The signs of $q(c_i)$ alternate.  If every gap maximum
were larger than $C_n$, then $q-q_*$ would have the same alternating signs
as $q$ at the $c_i$.  It would therefore have a zero between each consecutive
pair $c_i,c_{i+1}$, as well as the two zeros $\pm1$.  These are $n$ distinct
zeros, whereas the leading terms cancel and
$\deg(q-q_*)\le n-1$; moreover $q-q_*$ is nonzero at every $c_i$.
This contradiction selects a gap on which $|q|\le C_n$.  Scaling back proves
\emph{(9)}.

For sharpness, take
\[
 y_k=\frac{\cos((2k-1)\pi/(2n))}{r_n},
 \qquad 1\le k\le n.
\]
These are the zeros of $q_*$, have extremes $\pm1$, and every adjacent gap
contains a scaled Chebyshev extremum where $|q_*|=C_n$.  No smaller universal
constant can work.
\end{proof}

\begin{corollary}[collinear Erd\H{o}s case]
\label{cor:collinear-erdos-1041}
If the zero occurrences of a monic polynomial of degree $n\ge2$ lie on one line
in the open unit disc, two of them are joined by a curve of length strictly
below $2$ inside $\{|f|<1\}$.
\end{corollary}

\begin{proof}
Their diameter satisfies $D<2$.  Also $C_n\le1$ for $n\ge2$, with equality
only at $n=2$, so \emph{(9)} is strictly below one.  The selected segment has
length at most $D<2$.
\end{proof}

\paragraph{Formal boundary.}
\label{rem:sharp-collinear-formal-boundary}
Lean checks sign preservation, the alternating root-count mechanism, the
scaled Chebyshev endpoint and uniform bound, and
Theorem~\ref{prop:sharp-collinear-chebyshev-comparator}.  It does not check
the rigid normalization, the existence of the gap maxima, transport back to
the root line, or the equality-node locations used in
Theorem~\ref{thm:sharp-collinear-diameter}.  These are ordinary steps.  The
sharpness assertion concerns the maximum of $|f|$ on the selected adjacent-root
segment, and the displayed configurations attain it; no claim is made that
they exhaust the equality cases, and no claim is made that the selected
segment is a shortest path in $\{|f|<1\}$.

\subsection{The primitive sparse quintic}
\label{subsec:primitive-sparse-quintic}

The next selector is finite but unusually rigid.  For $0<r<2$ and
$0\le s_i\le1$, $x_i^2\le s_i$, put
\[
 E_i=s_i^4(s_i+r^2+2rx_i).
\]

\begin{theorem}[checked two-tail energy selector]
\label{prop:primitive-quintic-two-tail-energy-selector}
Suppose
\[
 \sum_{i=0}^4x_i=-r,\qquad
 \sum_{i=0}^4(2x_i^2-s_i)=r^2,\qquad
 \sum_{i=0}^4(4x_i^3-3s_ix_i)=-r^3.              \tag{10}
\]
Then at least one of the ten pairs $0\le i<j\le4$ satisfies
$E_i<1$ and $E_j<1$.
\end{theorem}

This is the exact finite conclusion selected as
\href{\repobase/ExternalVerification1041SolvedFamilies/Solution.lean\#L24}{the formal two-tail selector}.
It selects two distinct indices.  It does not assert that the corresponding
complex root values are distinct.

\begin{theorem}[primitive sparse quintic]
\label{thm:primitive-quintic-two-tail}
Let
\[
 p(z)=z^5+az^4+bz+c
\]
and suppose its five zero occurrences $w_0,\ldots,w_4$ lie in the closed unit
disc.  At least two distinct indices satisfy
\[
 |bw_i+c|\le1.                                    \tag{11}
\]
If $a\ne0$, two indices can be chosen with strict inequalities.  If $a=0$,
every index satisfies \emph{(11)}, and equality holds exactly when
$|w_i|=1$.

For open-disc zeros, two zero occurrences are joined inside $\{|p|<1\}$ by a
curve of length below $2$: use the two radial spokes through $0$ when their
values are distinct, and the constant path when the selected occurrences
have the same value.
\end{theorem}

\begin{proof}
Rotate so that $a=re^{i\phi}$ becomes $r\ge0$ and write
$z_i=e^{-i\phi}w_i$.  The missing $z^3$ and $z^2$ coefficients and Newton's
identities give
\[
 \sum z_i=-r,\qquad \sum z_i^2=r^2,\qquad
 \sum z_i^3=-r^3.                                 \tag{12}
\]
At a zero,
\[
 |bw_i+c|=|w_i|^4|w_i+a|.
\]
Thus, for $x_i=\Re z_i$ and $s_i=|z_i|^2$, the squared tail is exactly
$E_i$ and \emph{(12)} becomes \emph{(10)}.  If $r>0$, the third moment
also gives $r^3=|\sum_i z_i^3|\le5<8$, so $r<2$ as required by the
selector.

For completeness, the selector is driven by the harmonic extension
\[
 H_r(x,s)=(-r/2-x)(1-x)^2
 +(1-s)\left(1-\frac r4-\frac{3x}{4}\right).
\]
The three moments give the exact sum
\[
 \sum_{i=0}^4H_r(x_i,s_i)=5-2r,                  \tag{13}
\]
while $H_r\le4-2r$ throughout the unit disc.  If $E_r(x,s)\ge1$, then
\[
 5(1-s)\le2r(x+r/2).
\]
Put $d=x+r/2$, $u=1-x$, $A=1-r/4-3x/4$, and
$P=-u^2+(2r/5)A$.  The unsafe inequality gives $0\le d<2$.  Since
$H_r=-du^2+(1-s)A$, it is nonpositive when $A\le0$; when $A>0$,
the preceding bound gives $H_r\le dP$.  The remaining estimate is the exact
sum-of-squares identity
\[
 \frac1{31}-P
 =\frac25\left(d-\frac{38}{31}
               +\frac14(u-\frac3{31})\right)^2
  +\frac{31}{40}\left(u-\frac3{31}\right)^2\ge0,
\]
which gives $P\le1/31$, and hence $H_r(x,s)\le2/31$ whether $P$ is
positive or nonpositive.  Thus every unsafe tail has this smaller score.  If four tails were
unsafe, \emph{(13)} would give
\[
 5-2r\le4-2r+\frac8{31}<5-2r,
\]
a contradiction.  This proves the two-index conclusion.  The case $r=0$ is
the direct identity $|bw_i+c|=|w_i|^5$.

Finally, at a zero $w$,
\[
 p(tw)=(1-t)c+(t-t^4)(bw+c)-t^4(1-t)w^5.
\]
For $0\le t\le1$ the three nonnegative coefficients sum to $1-t^5$.
The tail bound, $|w|\le1$, and $|c|\le1$ therefore give the radial
unit-sublevel path.  Under the open-disc hypothesis the relevant bounds and
the total spoke length $|w_i|+|w_j|$ are strict.
\end{proof}

\paragraph{Formal boundary.}
\label{rem:primitive-quintic-formal-boundary}
Lean checks the five-index moment hypotheses, the harmonic cap, the
sum-of-squares unsafe estimate, and the ten-pair conclusion in
Theorem~\ref{prop:primitive-quintic-two-tail-energy-selector}.  The rotation,
derivation of Newton moments from the roots, root-tail identity, Abel
decomposition, path assembly, and any passage from distinct indices to
distinct complex values remain ordinary mathematics.

\subsection{Translated cubic quotient fibres}
\label{subsec:translated-cubic-quotient-fibres}

\begin{theorem}[checked cubic safe spoke]
\label{lem:cubic-safe-root-spoke}
If $r,s,v\in\mathbb C$ have modulus below one, at least one
$u\in\{r,s,v\}$ satisfies
\[
 \left|(tu-r)(tu-s)(tu-v)\right|\le1
 \qquad(0\le t\le1).
\]
\end{theorem}

This is exactly
\href{\repobase/ExternalVerification1041SolvedFamilies/Solution.lean\#L14}{the formal cubic safe-spoke theorem}.

\begin{theorem}[translated cubic quotient fibres]
\label{thm:translated-cubic-quotient-fibres}
Let $q\ge2$, $h\in\mathbb C$, $P$ be monic cubic, and
\[
 f(z)=P((z-h)^q).
\]
If every zero of $f$ lies in the open unit disc and $f$ has at least two
distinct zero values, then two zeros are joined through $h$ by a two-segment
path of length below $2$ inside $\{|f|<1\}$.  Equivalently this closes the
coefficient family
\[
 (z-h)^{3q}+A(z-h)^{2q}+B(z-h)^q+C
\]
in every degree $3q\ge6$.
\end{theorem}

\begin{proof}
Fix a $q$th root $y$ of a quotient root and a primitive $q$th root of unity
$\zeta$.  The full fibre consists of $h+y\zeta^k$.  Since every fibre point
lies in the open unit disc,
\[
 \frac1q\sum_{k=0}^{q-1}|h+y\zeta^k|^2
   =|h|^2+|y|^2<1.                                \tag{14}
\]
Thus $|y|<1$, and the corresponding quotient root has modulus
$|y|^q<1$.  This verifies the open-disc hypothesis for all quotient roots
before applying the following selector.
The identity retains more information than these separate bounds:
$|h|<1$ and every fibre radius is strictly below $\sqrt{1-|h|^2}$.
Consequently, any two safe fibre spokes constructed below have total length
strictly less than $2\sqrt{1-|h|^2}$.  The position of the symmetry centre
therefore controls the metric budget; the bound $2$ discards this information.

Write $P(w)=(w-r)(w-s)(w-v)$ and assign to $r$ the charge
$A_r=\Re(r\overline{s+v})$, cyclically.  The exact sum is
\[
 A_r+A_s+A_v=|r+s+v|^2-(|r|^2+|s|^2+|v|^2)>-3,
\]
so one charge, say $A_r$, exceeds $-1$.  For $0\le t\le1$, AM--GM and the
distance-square identity give
\[
 |tr-s|\,|tr-v|<1+t+t^2.
\]
Consequently, for $0\le t<1$,
\[
 |P(tr)|<(1-t)(1+t+t^2)=1-t^3\le1.              \tag{15}
\]
At $t=1$ the polynomial vanishes.  The entire spoke is therefore strictly
contained, including its origin endpoint.

For a nonzero safe quotient root from \emph{(15)}, two distinct
fibre points satisfy
\[
 f(h+ty\zeta^k)=P(t^qr),
\]
and their two spokes through $h$ have total length
$2|y|<2\sqrt{1-|h|^2}\le2$.  If the selected
quotient root is zero, choose a nonzero quotient root and use the direct
estimate $|P(ts)|<2t(1-t)\le1/2$ for $0<t<1$, with zero values
at both endpoints; if none exists, $f$ has only one distinct
zero value, contrary to the hypothesis.
\end{proof}

\paragraph{Formal boundary.}
\label{rem:cubic-quotient-formal-boundary}
Lean checks the charge identity and selection, the two-distance envelope,
the cubic product bound, the safe-spoke disjunction in
Theorem~\ref{lem:cubic-safe-root-spoke}, and an exact quartic falsifier for
this charge method.  It does not check the finite $q$-fibre construction,
root-of-unity average \emph{(14)}, zero-root fallback, pullback, or geometric
path assembly.  These remain ordinary proof steps.

\paragraph{Shared assurance boundary.}
\label{bdry:solved-polynomial-families}
The three Comparator declarations route respectively to
Theorems~\ref{prop:sharp-collinear-chebyshev-comparator},
\ref{prop:primitive-quintic-two-tail-energy-selector}, and
\ref{lem:cubic-safe-root-spoke}.  They support, but are strictly weaker than,
the ordinary assembled Theorems~\ref{thm:sharp-collinear-diameter},
\ref{thm:primitive-quintic-two-tail}, and
\ref{thm:translated-cubic-quotient-fibres}.  The displayed Lean endpoints and the assembled geometric theorems therefore
have different formalization boundaries.

\section{The current frontier: what survives near Fekete configurations}
\label{sec:frontier}

The August~29 public research update changes the useful first reading of this
problem.  It does not add a proof of Erd\H{o}s~\#1041; it identifies the
geometric regime in which the remaining difficulty is concentrated and kills
several attractive but false shortcuts.  The source of the update is pinned to
commit
\href{https://github.com/wcook04/plectis-erdos/tree/f214a6b45528dc5eefe20ffadc35f2e981627d4c/research_corpus/Erdos1041}{\texttt{f214a6b4}}
and the dated synthesis is
\href{https://github.com/wcook04/plectis-erdos/blob/f214a6b45528dc5eefe20ffadc35f2e981627d4c/research_corpus/Erdos1041/FRONTIER.md}{\texttt{FRONTIER.md}}.
The distinctions below are part of the result: a theorem, a refutation, a
certificate, and a measurement do not have interchangeable force.

\subsection*{Near-Fekete stability removes the radial distraction}

Write $a_i=\rho_i u_i$, with $|u_i|=1$, and put
$D=|\operatorname{disc}(f)|/n^n$.  The quantitative
Fekete--Hadamard estimate in
\href{https://github.com/wcook04/plectis-erdos/blob/f214a6b45528dc5eefe20ffadc35f2e981627d4c/research_corpus/Erdos1041/NearFeketeRadialAngularSplit.md}{\texttt{NearFeketeRadialAngularSplit.md}}
says that, when $D\ge1-\eta$ and
$\eta\le1/(80n^2)$,
\[
  1-\rho_i^2\le \frac{n\eta}{n-1},\qquad
  |a_i-a_j|\ge \frac{2-2\sqrt\eta-n\eta}{n-1},
\]
and the roots lie within $7\sqrt\eta$ of a rotated regular $n$-gon, after a
bijection.  The mechanism is a normalised Vandermonde Gram determinant:
near equality in Hadamard's inequality forces both the radii and the angular
separation to be near extremal.  This is an ordinary mathematical result in
the public research corpus, not a new Lean declaration.

The associated radial monotonicity theorem is the more important reader
conclusion.  Under the stronger near-Fekete condition
$\eta\le1/(10n^4)$ used by this corollary, replacing the radii by one can only
increase the values along the corresponding radial spoke.  Thus
radial deficits help containment; the unresolved near-Fekete problem is
angular.  On exact regular-gon directions this becomes unconditional for
$n\le6$: arbitrary radii in the stated band give
\[
 |f(s\,\omega\zeta^i)|\le1-s^n
 \quad(0\le s\le\rho_i),
\]
so every two-radii path is contained.  The good spoke in the general
on-circle averaging identity may move with $s$, however; that is why an
existence statement at one level does not supply a fixed pair of roots.

\subsection*{A bounded-radius concyclic class}

There is nevertheless a clean theorem for a different on-circle mechanism.
Let $f$ be monic of degree $n\ge3$ with distinct zeros on a circle of radius
$\rho$.  If $2\rho^n\le1$, two adjacent zeros are joined by their straight
chord, whose length is at most $2\rho\sin(\pi/n)<2$, and the chord lies in
$\{|f|<1\}$.  (If a zero is repeated, the short-connection conclusion is
immediate; the distinct case is the substantive one.)

After normalisation to
the unit circle, self-inversive realification and alternation against $z^n-c$
select a zero-free gap on which the polynomial modulus is at most $2$; a
harmonic normal-derivative argument transfers this arc bound strictly to the
chord.  The complete ordinary proof is
\href{https://github.com/wcook04/plectis-erdos/blob/f214a6b45528dc5eefe20ffadc35f2e981627d4c/research_corpus/Erdos1041/ConcyclicAlternation.md}{\texttt{ConcyclicAlternation.md}}.

The argument is an ordinary proof outside Lean.  The
\href{https://github.com/wcook04/plectis-erdos/blob/f214a6b45528dc5eefe20ffadc35f2e981627d4c/research_corpus/Erdos1041/scripts/check_erdos1041_concyclic_exact_witness.py}{exact-rational checker}
checks finitely many load-bearing identities and configurations, while the
\href{https://github.com/wcook04/plectis-erdos/blob/f214a6b45528dc5eefe20ffadc35f2e981627d4c/research_corpus/Erdos1041/scripts/check_erdos1041_concyclic_alternation.py}{numerical checker}
is regression and stress-test evidence.  The arc constant $2$ is sharp on the
regular $n$-gon, so this chord argument does not reach radii tending to $1$;
the unrestricted concyclic case and Erd\H{o}s~\#1041 remain open.

\subsection*{Cyclic trinomial fibres}

The spoke identity and strict bound are formalized in
\lref{Erdos1041/CyclicTrinomialFiberCase.lean}{46}{trinomialRoot_spoke_factorization}
and \lref{Erdos1041/CyclicTrinomialFiberCase.lean}{103}{trinomialRoot_spoke_norm_lt_one_of_norm_lt_one}.

There is also a positive structured family in every cyclic lift degree.  Fix
integers $1\le r\le m$ and $q\ge1$, and consider the translated polynomial
\[
 f(z)=(z-h)^{qm}+a(z-h)^{qr}+c.
\]
Writing $w=(z-h)^q$ reduces the root equation to
$w^m+aw^r+c=0$.  At such a quotient root the middle coefficient can be
eliminated exactly: for $0\le u\le1$,
\[
 u^mw^m+au^rw^r+c
   =(1-u^r)c-(u^r-u^m)w^m.
\]
The two coefficients on the right are nonnegative and have total at most one.
The \lword{Erdos1041/CyclicTrinomialFiberCase.lean}{103}{trinomialRoot_spoke_norm_lt_one_of_norm_lt_one}{formal strict-spoke theorem}
therefore proves that, if $|w|<1$ and $|c|<1$, the complete radial spoke
is strictly contained in the open unit lemniscate.  This is an all-root
statement; no Vieta-small root selection is needed for the containment bound.

For a nontrivial cyclic fibre, two selected displacements $y_1,y_2$ with
$|y_1|,|y_2|<1$ also satisfy $|y_1|+|y_2|<2$.  The formal source checks the
factorization, strict spoke estimate, and this metric budget.  It does not yet
check the ordinary finite argument selecting a suitable quotient root and two
distinct members of its fibre, lifting both quotient spokes, or assembling
their union as a path in the relevant component.  Thus the Lean theorem is the
load-bearing analytic kernel for this translated cyclic-trinomial family, not
the complete path theorem and not a proof of unrestricted Erd\H{o}s~\#1041.

\subsection*{Coefficient-controlled cyclic tetranomial spokes}

The first additional lower monomial admits an exact Abel-tail criterion.  Let
\[
  g(w)=w^m+aw^r+bw^s+c,\qquad m>r>s\ge1,
\]
assume all roots of $g$ lie in the open unit disk, and let $w_1,w_2$ be roots
of the two smallest moduli.  If
\[
  |c|+|b|\,|w_2|^s<1,
\]
then the complete radial spokes from $w_1$ and $w_2$ to the origin lie in
$\{|g|<1\}$, and their broken line has length strictly below $2$.  The simpler
coefficient condition
\[
  |b|+|c|\le1
\]
makes every root spoke safe; the coefficient $a$ is unrestricted.

The mechanism is visible in one identity.  At a root $w$ and for
$0\le u\le1$,
\[
\begin{split}
 g(uw)={}&(1-u^s)c
  -(u^s-u^r)(aw^r+w^m)
  -(u^r-u^m)w^m,\\
 aw^r+w^m={}&-(c+bw^s).
\end{split}
\]
The three scalar coefficients are nonnegative and sum to $1-u^m$.  Thus the
root-dependent budget $|c|+|b||w|^s<1$, together with $|w|<1$ and $|c|<1$,
puts every vector in the Abel decomposition strictly inside the unit ball.
Lean checks the \lword{Erdos1041/CyclicTetranomialCoefficientCase.lean}{27}{tetranomialRoot_spoke_factorization}{factorization}
and the resulting \lword{Erdos1041/CyclicTetranomialCoefficientCase.lean}{145}{tetranomialRoot_spoke_norm_lt_one_of_rootBudget}{strict spoke theorem} under
the exact weak exponent hypotheses $1\le s\le r\le m$; it also checks the
coefficient-only corollary above.

There is a stronger, genuinely two-index formal theorem.  Let $S$ index a
finite family of roots, $N=|S|\ge2$, and
\[
  M=\sum_{i\in S}w_i^s.
\]
The \lword{Erdos1041/TetranomialL2Selector.lean}{33}{sum_normSq_const_add_mul}{formal energy identity} gives
\[
 \sum_{i\in S}|c+bw_i^s|^2
 =N|c|^2+|b|^2\sum_{i\in S}|w_i^s|^2
   +2\mathop{\rm Re}(\overline c\,bM).
\]
Consequently, if every $|w_i|<1$, $|c|<1$, and
\[
 N\bigl(|b|^2+|c|^2\bigr)
   +2\mathop{\rm Re}(\overline c\,bM)<N-1,
\]
then two distinct indices $i,j\in S$ satisfy
$|c+bw_i^s|<1$ and $|c+bw_j^s|<1$.  Feeding those inequalities into the
\lword{Erdos1041/TetranomialL2Selector.lean}{208}{exists_two_tetranomialRoot_safeSpokes_of_moment_coeff_budget}{formal Abel-spoke consumer} proves that both complete root spokes lie
strictly in $\{|g|<1\}$.  The signed cross term is essential: this interface
can certify configurations outside the coefficient-only triangle
$|b|+|c|\le1$.  The formal hypotheses do not require $i\mapsto w_i$ to be
injective, so distinct indices need not denote distinct root values; the
ordinary two-root consequence requires a root enumeration without repetition.

For $q\ge2$ and $f(z)=g((z-h)^q)$, the ordinary regular-fibre mean-square
identity puts every quotient root inside the unit disk.  Lifting two selected
quotient spokes gives a path through $h$ whose two displacement lengths sum
to less than $2$.  The new Lean theorem assumes the finite root family, its
signed moment $M$, and the displayed budget; it does not derive that budget
from arbitrary tetranomial coefficients, prove that the supplied family is a
complete root multiset, perform cyclic-fibre lifting, or construct the final
path object.  Without the signed-moment, root-dependent, or coefficient-only
budget, the tetranomial case remains open; this family does not solve
unrestricted Erd\H{o}s~\#1041.

\subsection*{Translated quartic quotient fibres}\label{sec:translated-quartic-quotient-fibres}

Pendyala's degree-four theorem can also be lifted through every nontrivial
cyclic power.  Let $P$ be a monic quartic, let $q\ge2$, and set
\[
  f(z)=P((z-h)^q).
\]
If all listed zeros of $f$ lie in the open unit disk, averaging the squared
moduli over each complete $q$-point fibre shows that every quotient root
$w$ of $P$ satisfies $|w|<1$.  Pendyala's chord-or-radial proof then supplies
the quotient geometry.  The extra issue is metric: a short chord in the
$w$-plane need not lift isometrically through $y\mapsto y^q$.

Put $\alpha=1/q$.  Along a quotient chord whose supporting line has distance
$d$ from the origin, the root-lift density is controlled pointwise by
\[
  \bigl(\sqrt{d^2+x^2}\bigr)^{\alpha-1}
    \le x^{\alpha-1}\qquad(x>0),
\]
and Lean checks the \href{\repobase/research_corpus/Erdos1041/QuarticQuotientFiberCase.lean\#L37}{exact primitive}
\[
  \alpha\int_0^A x^{\alpha-1}\,dx=A^\alpha.
\]
Splitting at the perpendicular foot, and at the origin when the chord crosses
it, yields the ordinary power-map chord-lift estimate
\[
  \operatorname{length}(\widetilde{[a,b]})
    \le |a|^{1/q}+|b|^{1/q}.
\]
The \href{\repobase/research_corpus/Erdos1041/QuarticQuotientFiberCase.lean\#L59}{formal endpoint consumer} also proves that if
$0\le a,b<1$, $\alpha>0$, and a candidate length $L$ is at most
$a^\alpha+b^\alpha$, then $L<2$.  The two branches are handled separately, and
the distinction matters.  The chord-lift estimate applies to a single straight
quotient chord, and it bounds the lifted curve by the sum of the two endpoint
radii.  It does not bound the lift of an arbitrary broken line with a nonzero
interior vertex by that sum, so the radial branch is kept in its own form: the
fibre bound puts every quotient root inside a centred disc of radius strictly
below one, and Pendyala's four-point radial construction is applied there,
with the lifted spokes measured through the origin of the quotient plane.
With the two branches separated in that way both survive the power-map lift,
giving the asserted short path for every translated quartic quotient fibre, in
each degree $4q\ge8$.

The attribution and proof boundary are exact.  Pendyala proves the quartic
geometric theorem and its four-point radial lemma.  The local Lean module
checks the antitone density inequality, its integral, the strict powered
endpoint budget, and the final length fan-in.  It does not formalize
Pendyala's geometric lemma, the continuous covering-space construction of the
root lift, or the ordinary chord/radial case assembly.  This is a structured
all-scale family, not a proof of unrestricted Erd\H{o}s~\#1041.

\subsection*{A critical-value budget in every degree}

The critical values satisfy a sharp aggregate bound in every degree $n\ge2$.  This gives a budget for the merge levels; it does not
bound the lengths of the paths that reach them.

\begin{theorem}[critical-value budget in every degree]
\label{res:critical-value-budget}
Let $f$ be monic of degree $n\ge2$, with roots in a closed disc of radius
$R\ge0$.  If $c_1,\ldots,c_{n-1}$ are its critical points counted with
multiplicity, then
\begin{equation}\label{eq:critical-value-power-budget}
 \sum_{j=1}^{n-1}|f(c_j)|^{1/(n-1)}\le(n-1)R^{n/(n-1)}.
\end{equation}
Consequently
\[
 \sum_{j=1}^{n-1}|f(c_j)|^{1/n}\le(n-1)R.
\]
The constant is attained by $f(z)=(z-\tau)^n-\lambda$ with enclosing
disk centred at $\tau$ and radius $R=|\lambda|^{1/n}$.
\end{theorem}

The finite inequality behind the theorem concerns arbitrary points of the
disk, without asking them to arise as critical points.  Write $(\mathrm{FP}_m)$
for
\[
 \sum_{j=1}^m\left(\prod_{k=1}^m
       |1-\overline{c_j}c_k|\right)^{1/m}\le m,
 \qquad |c_j|\le1.
\]
The following pointwise bound explains why these finite inequalities
control critical values in every degree.

\begin{lemma}[reflected-derivative bound]\label{res:reflected-critical-value}
If $f$ is monic of degree $n\ge2$ with roots in the closed unit disk,
and $c_1,\ldots,c_{n-1}$ list its critical points with multiplicity, then
\begin{equation}\label{eq:critical-reflected-product}
 |f(c_j)|\le\prod_{k=1}^{n-1}|1-\overline{c_j}c_k|.
\end{equation}
\end{lemma}

\noindent\href{https://github.com/wcook04/plectis-erdos/blob/bb4e24651b37cd096d3749ef299c3317930b3a3b/ErdosProblems/Erdos1041/PaperReflectedCompletion.lean#L254}{The reflected-derivative inequality is checked in Lean, including closed-disc roots and critical-point multiplicities.}

\begin{proof}
First put all roots $a_i$ strictly inside the disk.  Gauss--Lucas does the
same for the $c_k$, so
\[
 N(z)=nf(z)-zf'(z),\qquad
 G(z)=n\prod_k(1-\overline{c_k}z)
\]
have $G\ne0$ on the closed disk.  On $|\zeta|=1$,
$|G(\zeta)|=|f'(\zeta)|$, while
\[
 \operatorname{Re}\frac{\zeta f'(\zeta)}{f(\zeta)}
 =\sum_i\left(\frac12+
       \frac{1-|a_i|^2}{2|1-a_i\bar\zeta|^2}\right)\ge\frac n2.
\]
Writing the logarithmic derivative as $w$, the identity
$|n-w|^2-|w|^2=n^2-2n\operatorname{Re}w\le0$ gives
$|N|\le|G|$ on the circle.  The maximum-modulus principle applied to
$N/G$ gives the same comparison inside.  At $c_j$, $N(c_j)=nf(c_j)$;
dividing by $n$ and conjugating each factor proves
\eqref{eq:critical-reflected-product}.
For closed-disk roots apply this result to
$f_r(z)=r^nf(z/r)$, $0<r<1$, whose critical points are $rc_k$:
\[
 r^n|f(c_j)|\le\prod_k|1-r^2\overline{c_j}c_k|.
\]
Let $r\uparrow1$.  This also handles boundary critical points and all
multiplicities without selecting local branches.
\end{proof}

Set $m=n-1$. By Gauss--Lucas and the lemma, $(\mathrm{FP}_m)$ gives
\[
 \sum_j |f(c_j)|^{1/m}\le m
\]
on the unit disk. Apply this to $R^{-n}f(h+Rz)$ for a root disk
centred at $h$ with $R>0$. Scaling the critical values back gives
\eqref{eq:critical-value-power-budget}. If $R=0$, all critical values
vanish. For the linear consequence put $y_j=|f(c_j)|^{1/m}$ in the
unit-disk normalization and use
\[
 y^{m/(m+1)}\le\frac{my+1}{m+1}\quad(y\ge0).
\]
Summation gives $\sum_j|f(c_j)|^{1/n}\le m$, and scaling back multiplies
this sum by $R$. The weighted Poisson proof below establishes $(\mathrm{FP}_m)$
for every $m\ge1$, including its equality case. Its quadratic form gives the
stronger bound
\[
 \sum_{j=1}^{n-1}|f(c_j)|^{2/(n-1)}\le(n-1)R^{2n/(n-1)}.
\]
The checked endpoints are the
\href{https://github.com/wcook04/plectis-erdos/blob/3f1a5e9b284b9c3348eebd549e24c2f3972606ae/ErdosProblems/Erdos1041/PaperCriticalValueMeanR10.lean#L101}{critical-value mean identity}
and the
\href{https://github.com/wcook04/plectis-erdos/blob/3f1a5e9b284b9c3348eebd549e24c2f3972606ae/ErdosProblems/Erdos1041/PaperCriticalValueMeanR10.lean#L121}{all-degree critical-value power budget};
the
\href{https://github.com/wcook04/plectis-erdos/blob/3f1a5e9b284b9c3348eebd549e24c2f3972606ae/ErdosProblems/Erdos1041/PaperCriticalValueMeanR10.lean#L85}{critical-disc moment bound}
gives all moments $0<t\le2/(n-1)$. Their source-bound compilation and named
axiom checks pass.

The stronger exponent controls concentration of the critical values.
For $R>0$, write $r_j=|f(c_j)|^{1/n}$ and $p=n/(n-1)$. Then
\[
 \sum_j(r_j/R)^p\le n-1,\qquad
 \#\{j:r_j\ge tR\}\le\frac{n-1}{t^p}\quad(t>0).
\]
The counting bound follows by retaining just those summands. Thus a large
critical value consumes more of the budget than the linear estimate records;
the radial equality family has every $r_j=R$. This distributional control
still leaves the geometry of the joining paths to be supplied.

\paragraph{The weighted Poisson proof and the small-cardinality conclusions.}
Let $w_j\ge0$, $\sum_jw_j=1$, and $|c_j|\le1$, with repetitions allowed.
Write $G(z)=\prod_k|1-\overline{c_k}z|^{w_k}$, taking every zero exponent
factor to be one. For strictly interior centres choose analytic logarithms
zero at the origin and put
\[
 g(z)=\exp\sum_jw_j\log(1-\overline{c_j}z)
     =1+\sum_{\nu\ge1}a_\nu z^\nu,\qquad
 P=\sum_jw_jP_{c_j}.
\]
On the unit circle $P=1-2\Re(\zeta g'/g)$. Weighted Poisson majorisation,
followed by absolutely justified circle integration, gives
\[
 \sum_jw_jG(c_j)^2\le\int |g|^2P\,dm
 =1-\sum_{\nu\ge1}(2\nu-1)|a_\nu|^2.
\]
All functions are analytic on a neighbourhood of the closed disc in this
interior case. For closed-disc centres replace $c_j$ by $rc_j$, $0<r<1$.
The actual coefficient of degree $\nu$ becomes $r^\nu a_\nu$, where the
fixed coefficients are defined from the analytic function near zero.
Pass to $r\uparrow1$ first with an arbitrary finite coefficient sum.
Nonnegative finite deficits then give summability and the infinite inequality;
no boundary holomorphic logarithm is required. Lean checks the
\href{https://github.com/wcook04/plectis-erdos/blob/3f1a5e9b284b9c3348eebd549e24c2f3972606ae/ErdosProblems/Erdos1041/WeightedCauchyTransportR10.lean\#L66}{exact radial transport of each Cauchy coefficient},
the
\href{https://github.com/wcook04/plectis-erdos/blob/3f1a5e9b284b9c3348eebd549e24c2f3972606ae/ErdosProblems/Erdos1041/WeightedBoundaryDeficitR10.lean\#L108}{finite boundary-deficit inequality},
and the
\href{https://github.com/wcook04/plectis-erdos/blob/3f1a5e9b284b9c3348eebd549e24c2f3972606ae/ErdosProblems/Erdos1041/WeightedBoundaryDeficitR10.lean\#L153}{summable full boundary deficit}.

If equality holds, all positive-degree coefficients vanish, so $g=1$ near
zero and $\sum_jw_j\overline{c_j}/(1-\overline{c_j}z)=0$ there.
Group equal centres before clearing denominators. A nonzero centre of positive
total weight gives a nonzero pole coefficient, which is impossible. Thus
equality is equivalent to $c_j=0$ on positive support. In particular,
strictly positive weights force all centres zero; $(w,c)=((1,0),(0,1))$
shows why this wording cannot be extended to zero weights.
The exact grouped-rigidity source is preserved, not replaced by an argument
requiring distinct original centres.

For $M=\sum_jw_jG(c_j)$ the identity
\[
 M^2+\sum_jw_j(G(c_j)-M)^2=\sum_jw_jG(c_j)^2
\]
gives weighted Cauchy--Schwarz, with equality precisely when the observations
are constant on positive support. Equal weights therefore give
$(\mathrm{FP}_m)$ for every $m\ge1$, with equality only at the zero
configuration. The
\href{https://github.com/wcook04/plectis-erdos/blob/3f1a5e9b284b9c3348eebd549e24c2f3972606ae/ErdosProblems/Erdos1041/PaperWeightedCoverageR11.lean\#L138}{three-point theorem}
and
\href{https://github.com/wcook04/plectis-erdos/blob/3f1a5e9b284b9c3348eebd549e24c2f3972606ae/ErdosProblems/Erdos1041/PaperWeightedCoverageR11.lean\#L146}{four-point theorem}
state these conclusions with their equality cases.

The previous three-point H\"older proof is withdrawn as an alternative proof
from this record. Its
\href{https://github.com/wcook04/plectis-erdos/blob/3f1a5e9b284b9c3348eebd549e24c2f3972606ae/ErdosProblems/Erdos1041/ThreePointDefectR11.lean\#L16}{scalar defect identity}
survives independently, but that algebra alone is not credited with
the omitted H\"older step or the whole alternative argument. The four-point
matching and numerical split are likewise not used. The present argument
replaces those passages; it does not claim to formalise their mechanisms.

\paragraph{Poisson majorization and termwise integration.}
For $g$ holomorphic on the unit disc and continuous on its closure, put
$P_c(\zeta)=(1-|c|^2)/|\zeta-c|^2$ for $|\zeta|=1$.  If $|c|<1$,
\[
 |g(c)|^2\le\frac1{2\pi}\int_0^{2\pi}
 P_c(e^{it})|g(e^{it})|^2\,dt.
\]
The \href{https://github.com/wcook04/plectis-erdos/blob/d4fed71423840f70f10edf27b9ad27c22fc4f49a/ErdosProblems/Erdos1041/PoissonNormSquare.lean#L114}{norm-square majorization}
uses the Poisson formula for a real part and the nonnegativity of a square.
The \href{https://github.com/wcook04/plectis-erdos/blob/d4fed71423840f70f10edf27b9ad27c22fc4f49a/ErdosProblems/Erdos1041/PoissonKernelBridge.lean#L27}{weighted kernel identity}
identifies the pointwise Poisson mixture.  Separately,
\href{https://github.com/wcook04/plectis-erdos/blob/d4fed71423840f70f10edf27b9ad27c22fc4f49a/ErdosProblems/Erdos1041/CircleSeriesTransport.lean#L49}{absolute series transport}
justifies termwise circle integration when the terms are circle integrable
and admit a summable uniform norm bound; its Taylor double-product form
assumes absolute summability of both coefficient sequences.
The
\href{https://github.com/wcook04/plectis-erdos/blob/3f1a5e9b284b9c3348eebd549e24c2f3972606ae/ErdosProblems/Erdos1041/PaperAnalyticR10.lean\#L1}{checked analytic assembly}
is accompanied by an
\href{https://github.com/wcook04/plectis-erdos/blob/3f1a5e9b284b9c3348eebd549e24c2f3972606ae/ErdosProblems/Erdos1041/PaperAnalyticR10Audit.lean\#L1}{80-declaration audit}.
The repaired R10 dependency chain and the
\href{https://github.com/wcook04/plectis-erdos/blob/3f1a5e9b284b9c3348eebd549e24c2f3972606ae/ErdosProblems/Erdos1041/PaperAnalyticR11.lean\#L1}{R11 aggregate}
and
\href{https://github.com/wcook04/plectis-erdos/blob/3f1a5e9b284b9c3348eebd549e24c2f3972606ae/ErdosProblems/Erdos1041/PaperAnalyticSmokeR11.lean\#L1}{smoke target}
pass. Across the R10, R11, and
\href{https://github.com/wcook04/plectis-erdos/blob/3f1a5e9b284b9c3348eebd549e24c2f3972606ae/ErdosProblems/Erdos1041/PaperAnalyticClosureAuditR11.lean\#L1}{closure audits},
320 named prints report only \path{propext}, \path{Classical.choice}, and
\path{Quot.sound}.

\paragraph{A uniform central region in every degree.}
For every $m\ge1$, the
\href{https://github.com/wcook04/plectis-erdos/blob/3be82b1a7340284aea72e9a5c8493cb020843921/ErdosProblems/Erdos1041/FreePointCentralCompletion.lean#L39}{central-region theorem with its analytic inputs discharged}
proves
\[
 |c_j|\le\sqrt{1-e^{-2}}\quad(1\le j\le m)
 \quad\Longrightarrow\quad
 \sum_{j=1}^m\left(\prod_{k=1}^m|1-\overline c_jc_k|\right)^{1/m}\le m.
\]
To see the mechanism, set
$h_j=m^{-1}\sum_k\log|1-\overline c_jc_k|$ and
$D_j=-\log(1-|c_j|^2)$.  The logarithmic kernel gives
$\sum_jh_j\le0$ and
$h_j^2\le(-m^{-1}\sum_i h_i)D_j$.  The radius bound gives $D_j\le2$.
Writing $r^2=-m^{-1}\sum_i h_i$, the resulting bounds
$\sum_jh_j=-mr^2$ and $|h_j|\le\sqrt2r$ feed the exponential energy
inequality, yielding $\sum_j e^{h_j}\le m$.  Thus the series estimates are
proved for the actual configuration, rather than assumed as suppliers.
This checks a central equal-weight case; the short paper's ordinary
all-disc weighted inequality has greater mathematical scope.

The earlier adaptive certificate remains useful for other configurations.

For $m\ge1$, let $|c_i|^2\le a<1$, with $a\ge0$, and put
$L=-\log(1-a)$. The
\href{https://github.com/wcook04/plectis-erdos/blob/457e0c74b2666f7fb2d825c967e49e439a600b23/ErdosProblems/Erdos1041/FreePointUniformRadius.lean#L84}{uniform-radius theorem}
gives the actual geometric-mean inequality
\[
 e^L\le1+2L\quad\Longrightarrow\quad
 \sum_{i=1}^m\left(\prod_{j=1}^m|1-\overline c_i c_j|\right)^{1/m}\le m.
\]
Its mechanism is an adaptive logarithmic certificate. Write
$H_i=m^{-1}\sum_j\log|1-\overline c_i c_j|$ and
$D_i=-\log(1-|c_i|^2)$. For nonnegative caps $M_i\ge H_i$, the
\href{https://github.com/wcook04/plectis-erdos/blob/0d34630e1cc9d2b1ac6edfa7cfdba83b31bdc8fe/ErdosProblems/Erdos1041/LogKernelCentralCertificate.lean#L99}{checked adaptive theorem}
requires only $\sum_i\Phi(M_i)D_i\le m$, where
$\Phi(t)=\int_0^1(1-s)e^{st}\,ds$.
The uniform radius gives $H_i,D_i\le L$, so $\Phi(L)L\le1$
suffices; the displayed exponential condition is exactly this inequality
when $L>0$, and $L=0$ is immediate.

\paragraph{Four points without a matching or numerical split.}
The preceding weighted proof with $m=4$ and $w_j=1/4$ gives
\[
 \sum_{j=1}^4\left(\prod_{k=1}^4|1-\overline{c_k}c_j|\right)^{1/4}\le4,
 \qquad |c_j|\le1,
\]
with equality if and only if all four centres vanish. The
\href{https://github.com/wcook04/plectis-erdos/blob/3f1a5e9b284b9c3348eebd549e24c2f3972606ae/ErdosProblems/Erdos1041/PaperWeightedCoverageR11.lean\#L146}{checked four-point theorem}
states both conclusions, and the
\href{https://github.com/wcook04/plectis-erdos/blob/3f1a5e9b284b9c3348eebd549e24c2f3972606ae/ErdosProblems/Erdos1041/PaperWeightedCoverageR11.lean\#L128}{strict form}
retains strictness as soon as a centre
is nonzero, including repeated centres and unit-circle centres.

The former $21/25$ central/outer split, its perfect-matching H\"older step,
and its polynomial, Taylor, logarithmic and radical-calculus certification
claims are removed from this record. A scalar stationary-point inequality
alone does not establish that whole route. The central-region and adaptive
proofs immediately above remain independent, restricted results, not
certificates for the removed four-point estimates. The replacement weighted
proof and its explicit three- and four-point endpoints pass their source-bound
compilation and named axiom checks.

The budget guarantees a critical value at most $R^n$, but does not force
the smaller quintic threshold $1/M_5$ below.  On the unit radial family all
critical values have modulus one, so a uniformly strict improvement is
unavailable.  Nor does the value budget bound inverse-ray lengths: a path
selection or metric comparison is still needed.

\subsection*{The degree-five target is now explicit}

The degree-five analysis sharpens the two-nearest-spoke threshold.  If a
critical point $c$ satisfies $|f(c)|\le1/M_n$, then both straight segments to
the two nearest roots remain in $\{|f|\le1\}$ and their total length is at
most $2$.  The envelope is obtained by maximising the exact product along the
second spoke; at degree five,
\[
 M_5=(1-t_*)(1+t_*)^3\sqrt{16t_*^2-4t_*+1},
 \qquad
 t_*={5\over16}+{3\sqrt{105}\over80},
 \qquad
 {1\over M_5}=0.2760461\ldots.
\]
This is an 11-fold improvement over the earlier deep-low threshold, but it
does not assert that every quintic has a critical value below this level.

The exact remaining degree-five statement is therefore not ``improve the
constant'' but the following selection problem:
\begin{quote}
For every monic quintic with roots in the closed unit disc, there is a
critical point $c$ with $|f(c)|\le1$ and two roots $a,b$ such that
$|f|\le1$ on $[c,a]\cup[c,b]$ and
$|c-a|+|c-b|\le2$.
\end{quote}
The public file calls this (SPOKE-5).  Together with the existing
two-segment reduction it would settle degree five, but (SPOKE-5) is not
proved.  Coverage experiments put the residual in rapid, nearly simultaneous
merging near the regular pentagon: the generic sampled families were covered,
whereas the surviving band has every named merge threshold active and
$D$ near $1$.  Those percentages locate the work; they are not a theorem.

At the same degree, the terminal connected component admits a sharper area
identity.  If $K_t=\{|f|\le t\}$ is connected and the centred exterior map is
\[
 \psi(\zeta)=t^{1/n}\zeta+a_0+
 \sum_{k\ge1}a_k\zeta^{-k},
\]
then Gr\"onwall's identity gives
\[
 \operatorname{Area}(K_t)=\pi\left(t^{2/n}-
 \sum_{k\ge1}k|a_k|^2\right),
 \qquad
 a_1=-\frac{c_{n-2}}{n\,t^{1/n}}
 \quad\text{after centring}.
\]
It replaces a P\'olya upper bound at that terminal node, but extending it to
proper components and connecting it to the parent theorem remains open.  The
degree-five statements and their evidence classes are recorded in
\href{https://github.com/wcook04/plectis-erdos/blob/f214a6b45528dc5eefe20ffadc35f2e981627d4c/research_corpus/Erdos1041/Degree5AssemblyAndSharpenedCuts.md}{\texttt{Degree5AssemblyAndSharpenedCuts.md}}.

\subsection*{The no-go boundaries force hub selection}

The same frontier is useful because it closes off three misleading readings.
First, the strict-argmin first-merge hub need not have two-arm length below
$2$.  A degree-five tie-locus witness, moved into the open unit disc, has
\[
 L(c_*)=2.0573432753\ldots>2,
\]
with a 50-digit two-instrument reconstruction and a margin about 97 times
that of the earlier degree-four witness.  This is computational certificate
evidence, not an exact rational theorem; the parent remains viable because a
different hub at the same configuration supplies the relevant pair.

Second, the aggregate estimate
$\sum_cL(c)\le2(n-1)R$ is refuted at degree four on an open violating region
and at one degree-five witness.  Its algebraic factor
\[
 \sum_{k=1}^{n-1}|f(c_k)|^{1/n}\le(n-1)R
\]
is proved for every $n\ge2$ in Theorem~\ref{res:critical-value-budget}.
The metric factor is not automatic in any of these degrees: the all-degree
value estimate supplies no path-length aggregate.
These two boundaries, including the distinction between the surviving
algebraic budget and the failed length aggregate, are in
\href{https://github.com/wcook04/plectis-erdos/blob/f214a6b45528dc5eefe20ffadc35f2e981627d4c/research_corpus/Erdos1041/MinimalHubArmBudgetRefutation.md}{\texttt{MinimalHubArmBudgetRefutation.md}}
and
\href{https://github.com/wcook04/plectis-erdos/blob/f214a6b45528dc5eefe20ffadc35f2e981627d4c/research_corpus/Erdos1041/SeparatrixAggregateReduction.md}{\texttt{SeparatrixAggregateReduction.md}}.

Third, the origin is not a uniform near-Fekete hub.  An exact rational
quintic with roots scaled by $999999/1000000$ has four of five origin spokes
escaping $\{|f|\le1\}$ at explicit rational sample points, even though
$1-D=3.19991\cdot10^{-4}$.  This witness rules out the origin-spoke selector
on every defect range that contains its value.  Ruling the selector out on
every positive neighbourhood of the Fekete locus would need a sequence of
witnesses with defect tending to zero, or a limiting argument proving that
assertion; neither is supplied here.  The mechanism is angular: the first
several Fourier modes can point against different spokes.  The certificate is
exact for the displayed witness; the claim that the phenomenon persists
throughout a limiting family is computational.  The witness was engineered
against the origin selector, and a different admissible hub at the same
configuration supplies a contained pair.  What fails here is the origin
selector; the admissible hubs are not exhausted.  The full witness and replay route are in
\href{https://github.com/wcook04/plectis-erdos/blob/f214a6b45528dc5eefe20ffadc35f2e981627d4c/research_corpus/Erdos1041/NearFeketeRadialAngularSplit.md}{\texttt{NearFeketeRadialAngularSplit.md}}.

\subsection*{The surviving carrier and the actual open endpoint}

These refutations leave a selector as the surviving carrier.  On the
ray-separated dense class, the public frontier records the canonical open row
\[
  \min_{c\ \mathrm{admissible}} L(c)\le2,
\]
where admissibility includes the two inverse-ray arms being contained in the
relevant lemniscate component.  Its measured values remain below $2$ at the
refuting configurations, and the existing attachment and lower
semicontinuity reductions would turn this row into the parent theorem.  No
proof of the row is currently available.  In particular, the nonzero hub
that rescues the origin-spoke witnesses is evidence for the selector, not a
universal construction.

The honest current endpoint is therefore a containment selector with three
separate obligations: choose the hub, keep both arms inside the component,
and control the metric fan-in.  The Lean module cited elsewhere in this note
checks the Newton value equation, ray-collision avoidance, and root-retention
inputs; the new frontier files above are ordinary proof, exact-certificate,
and computational research records, not formal authority for (SPOKE-5), the
no-go generalisations, or the parent theorem.  Keeping those authority
boundaries visible is part of making the frontier reusable.

\section{The Newton value equation}\label{sec:newton}

Away from the critical set, define the complex Newton field
\[
  N(z)=-\frac{f(z)}{f'(z)} .
\]
Let $I\subseteq\mathbb R$ be an interval. On the interval under consideration,
assume $z$ is differentiable, $f'(z(t))\ne0$, and $z'(t)=N(z(t))$;
put $w(t)=f(z(t))$.

\begin{theorem}[value equation]\label{res:value}
$w'(t)=-w(t)$ on $I$.
\end{theorem}

\noindent The computation is one line: $w'=f'(z)\,z'=f'(z)\cdot(-f(z)/f'(z))=-f(z)=-w$.
The kernel checks it as \lref{Erdos1041/NewtonFlowRaySeparation.lean}{50}{newtonFlow_value_hasDerivAt}, together with the
differential form of the first integral,
\lref{Erdos1041/NewtonFlowRaySeparation.lean}{64}{newtonFlow_scaledValue_hasDerivAt_zero}:
\[
  \frac{d}{dt}\Bigl(e^{t}f(z(t))\Bigr)=0,
  \qquad\text{equivalently}\qquad
  f(z(t))=e^{-(t-t_0)}f(z(t_0)),\qquad t_0\le t,\quad t_0,t\in I .
\]
For an existing trajectory with nonzero initial value, its value moves
inward on one positive ray; a zero value has no argument and remains zero.
Thus $|f|<1$ is preserved for later times in that trajectory's existing
interval. No global existence or description of a whole ray preimage follows
from this scalar equation. The real-time candidate endpoint is
\path{newtonFlow_real_value_decay} in \path{NewtonFlowTrajectory.lean}.

\begin{corollary}[ray separation]\label{res:ray}
Let $a<b$ and let the value trajectory $t\mapsto f(z(t))$ be continuous on
$[a,b]$. Assume the Newton equation and $f'(z(t))\ne0$ on $(a,b)$ only.
Then
\[
 f(z(b))=e^{a-b}f(z(a)).
\]
If these endpoint values are nonzero, they lie on one positive ray. Therefore
critical points with values on distinct positive rays cannot be endpoints
of such a finite connection. The trajectory in the $z$ plane need not be radial.
\end{corollary}

\noindent The interior identity passes to the endpoints by continuity, not
by evaluating $-f/f'$ at a critical endpoint. The candidate declarations
\path{PaperNewtonEndpoints.value_decay_at_continuous_endpoints} and
\path{PaperNewtonEndpoints.no_finite_connection_of_distinct_value_rays}
in \path{ErdosProblems/Erdos1041/PaperNewtonEndpoints.lean} make these
premises explicit. Their compilation and axiom checks are unrun in this
return. They construct no trajectories and supply no global monodromy theorem.

\section{Arguments, not moduli}\label{sec:arguments}

It is tempting to arrange a generic perturbation so that the critical values
are pairwise distinct, or that their moduli are pairwise distinct, and to
conclude that saddle connections are excluded.  Neither is enough.

Two distinct critical values can lie on one ray, and two critical values with
distinct moduli certainly can: the ray records the argument, and the modulus is
exactly the coordinate the flow contracts.  By Corollary~\ref{res:ray} the
invariant that excludes connections is the argument.  What a perturbation must
therefore achieve is pairwise distinct critical-value \emph{arguments}, which is
a condition on $n-1$ points modulo the circle.  Their positions in the plane
are unconstrained by it.

The cost of that condition is also checked, and it is small.

\begin{theorem}[ray-collision locus]\label{res:locus}
Let $a\ne b$ be complex.  Every common translation $\beta$ for which $a+\beta$
and $b+\beta$ lie on the same positive ray has the form
\[
  \beta=\frac{ra-b}{1-r},
  \qquad r\in\mathbb{R}_{>0},\ r\ne1 .
\]
\end{theorem}

\noindent So each pair of critical values contributes a one-real-parameter
forbidden locus in the translation plane, given in closed form
(\lref{Erdos1041/NewtonFlowRaySeparation.lean}{107}{translated_samePositiveRay_parameterization}).  A finite union of such
loci has empty interior, which is the shape one wants for an avoidance
argument.  Turning that into a perturbation of $f$ is not immediate: the
translation model must be replaced by an actual perturbation of the roots that
keeps them inside $\dbar$ and preserves the length slack.  That is the first
open producer of \S\ref{sec:open}.

\section{A proof gap in the unrestricted argument}\label{sec:gap}

The March manuscript's Proposition~12 claims the following load-bearing
statement.  For $u=-\log|f|$, a connected component $V$ of $\{u>c\}$ carrying
$m\ge2$ simple zeros, and the stated regularity and Morse hypotheses, it
constructs, for every $\varepsilon>0$, an embedded spanning tree
$G_\varepsilon\subset V$ with
\[
 \operatorname{len}(G_\varepsilon)
 \le\frac1{2\pi}\int_{2\alpha}^{\infty}P_V(t)\,dt+\varepsilon.
\tag{5.1}\label{eq:prop12-bound}
\]
The final theorem uses this estimate, so the issue below cannot be bypassed by
calling the proposition auxiliary.

At an interior index-one critical point $p$, the proof invokes a Morse chart
\[
 u=\mu+x^2-y^2
\]
and replaces the saddle by a three-ended neighbourhood having one connected
lower cross-section and two connected upper cross-sections.  That local model
is false as written.  Because $p\in V$ and $V$ is open, a sufficiently small
closed disc around $p$ lies entirely in $V$.  In that disc the full Morse chart
has four sectors: two components of $u>\mu$ and two components of $u<\mu$.
A global component argument cannot delete one local sector from a disc already
contained in $V$.

There is also a self-contained obstruction to the displayed unrestricted
budget.  For $f_a(z)=z^2-a^2$, $a=9/10$, let
\[
 B_R=\frac1{2\pi}\int_{|f_a(z)|<R}\left|\frac{f'_a(z)}{f_a(z)}\right|\,dA(z).
\]
Changing variables through both inverse branches gives
$B_R\le4(\sqrt{a^2+R}-a)$.  At $R=1$ this is less than
$223/125<9/5$, whereas every connected rectifiable tree containing the roots
$\pm a$ has length at least $9/5$.  Coarea identifies $B_1$ with the full
level-length budget, so no nonnegative cutoff can repair the estimate without
an added attachment cost.  This does not refute Erd\H{o}s~\#1041: the segment
$[-a,a]$ itself lies in $\{|f_a|<1\}$ and has length $9/5<2$.

Applying this obstruction to a numbered historical proposition requires its
complete archived hypotheses, which were not recovered in the return.  A
repair might cut an adjoining regular annulus along a separatrix or regular
flow arc before forming the block, retain a four-pronged saddle neighbourhood
and change the assembly, or replace the local construction by the ray-cut
decomposition proposed below.  Any repair must pay the missing attachment
cost or prove a different metric estimate.
The shorter descriptions of the same three-ended block do not repair the
four-sector topology.

Corollary~\ref{res:ray} supplies one independent input for a different route:
distinct critical-value arguments exclude saddle-to-saddle Newton
connections.  It does not itself prove the compact planar decomposition,
classify all orbit endpoints or provide the metric gluing estimate.  Those are
separate problems below.

\section{Finite evidence}\label{sec:finite}

A search was run over random monic polynomials with roots in the unit disc.  For
each sample the region $\{|f|<1\}$ was rasterised and shortest grid paths were
computed between every pair of roots.  The best upper bounds obtained were
\[
  \begin{array}{lrr}
    \text{degree } 5, & 500 \text{ trials}: & 1.1052648928\\
    \text{degree } 6, & 1500 \text{ trials}: & 0.8450414343\\
    \text{degree } 8, & 1500 \text{ trials}: & 0.6203916714\\
    \text{degree } 10, & 1500 \text{ trials}: & 0.4303640486 .
  \end{array}
\]
No counterexample candidate was found, and the measured values sit well below
the threshold $2$.  These are numerical candidate connectors.  A polygon whose
vertices satisfy $|f|<1$ need not lie in the open lemniscate, since an edge can
cross the boundary between two safe vertices, and root snapping adds further
edges of the same kind.  Turning a candidate into a proof needs a continuous
certificate on every edge, for instance strict positivity of the real
polynomial $1-|f(z(t))|^2$ on $[0,1]$ for each straight edge $z(t)=u+t(v-u)$,
established by exact coefficients or outward-rounded interval arithmetic with
subdivision, together with exact root enclosures.  No such certificate is
attached to the numbers above.  They are grid distances for the sampled
configurations, and the apparent decrease with degree is a property of the
sample, from which we draw no conjecture.  Even a finite collection of proved
instances would leave the universal statement open.
Future searches should target four named failure modes: near-degenerate
saddles, thin necks, boundary-critical configurations, and almost-connected
separatrices.  They should report those diagnostics.  Raster paths remain candidate
finders, never continuous certificates.

\section{Complements and further questions}\label{sec:open}

The dependency chain has five separate gates.  A proof, a minimally corrected
hypothesis set, or an explicit polynomial or planar counterexample is a
decisive answer to any one of them.

\subsection*{1. Repair or refute the saddle block}

\begin{problem}[corrected local saddle assembly]\label{prob:saddle1041}
Under the hypotheses of Proposition~12, replace the invalid three-ended local
model by a four-pronged block or by a block formed after a specified annular
cut.  Prove that for every $\eta>0$ its connector has total Euclidean length
below $\eta$, uniformly over the selected attachment points, and that the
corrected blocks assemble to an embedded tree satisfying
\[
 \operatorname{len}(G_\varepsilon)
 \le\frac1{2\pi}\int_{2\alpha}^{\infty}P_V(t)\,dt+\varepsilon;
\]
or give a polynomial or harmonic planar counterexample to that statement.
\end{problem}

The full-disc model $x^2-y^2$, an annulus in which lower branches rejoin, two
saddles joined by a separatrix and simultaneous saddle levels are mandatory
tests.  Repeating the one-lower/two-upper assertion does not answer the
problem.

\subsection*{2. The compact ray-cut decomposition}

Let $p$ have simple roots, simple nonzero critical points and let
$u=-\log|p|$.  For regular values $c<T$, take
\[
 M_{c,T}=\overline{V\cap\{c<u<T\}},
\]
and assume explicitly that this is a compact genus-zero surface, its lower
boundary is one smooth Jordan curve, its upper boundary consists of $m$ smooth
root curves, all interior critical points are nondegenerate index-one saddles,
the normalised gradient field
\[
 X=\frac{\nabla u}{|\nabla u|^2}=-\frac{p}{p'}
\]
is transverse to the level boundaries, and no maximal $X$-trajectory has two
saddle endpoints.

\begin{problem}[finite strip decomposition after ray cuts]\label{prob:reeb1041}
For every $\eta>0$, construct disjoint Morse neighbourhoods $N_j$ with
$\sum_j\operatorname{diam}N_j<\eta$ and a finite set of complete separatrix or
regular-flow cuts so that every component of the complement is flow-diffeomorphic
to a rectangle
\[
 [a_S,b_S]\times[0,1],\qquad u(\Phi_S(t,s))=t,
\]
with connected level sections and no uncut annular component.  Give an explicit
finite bound $E(s,m)$ for the number of strips, or exhibit the minimal missing
hypothesis or a counterexample.
\end{problem}

No particular formula such as $2s+1$ is presumed.  Boundary tangencies,
simultaneous levels, branch reunion through an annulus and non-Hausdorff orbit
spaces must all be handled explicitly.

\subsection*{3. Metric fan-in without losing the coefficient}

For a strip $S$, write
\[
 \Gamma_t^S=S\cap\{u=t\},\qquad
 P_S(t)=\mathcal H^1(\Gamma_t^S),
\]
and let $k_S$ be its number of root ends.  Write
\[
 \Phi_S(t)=\int_{\Gamma_t^S}|\nabla u|\,ds
\]
for its transverse flux.  Normalising the transverse measure by $\Phi_S$ gives
the average trajectory estimate
\[
 \int_{\Gamma_{t_0}^S}\operatorname{len}(\gamma_x)\,d\mu_{t_0}(x)
 =
 \frac1{\Phi_S(t_0)}\int_{a_S}^{b_S}P_S(t)\,dt .
\]
The flux of a strip is not an integer multiple of $2\pi$.  For $f(z)=z$ and
$u=-\log|z|$, an annular sector of angular width $\theta$ has $|\nabla u|=1/r$
and $ds=r\,d\phi$ on a circular level arc, so $\Phi_S=\theta$; for $z^n$ the
same computation gives $n\theta$.  A closed level curve enclosing roots does
carry an integer winding flux, and cutting a regular annulus destroys that.
Any recovery of the global coefficient $1/(2\pi)$ from per-strip estimates has
therefore to account for how the individual $\Phi_S$ sum, without duplication
and without a multiplicative loss.

\begin{problem}[additive-error strip gluing]\label{prob:metric1041}
Assuming Problem~\ref{prob:reeb1041}, select trajectories in all strips and
connect them through the saddle and root neighbourhoods so that, for every
$\eta>0$, the resulting embedded tree contains all $m$ roots and obeys
\[
 \operatorname{len}(G)
 \le\frac1{2\pi}\int_{2\alpha}^{\infty}P_V(t)\,dt+\eta.
\tag{6.1}\label{eq:metric-fanin1041}
\]
The total saddle, annular-cut and root-cap cost must be below $\eta$ without a
multiplicative loss in $1/(2\pi)$.  The coefficient $1/(2\pi)$
in~\eqref{eq:metric-fanin1041} is an open target and depends on a valid global
allocation of the strip fluxes $\Phi_S$; the remaining work is not a small
local cap added to an otherwise complete estimate.
\end{problem}

For the final strict inequality, use the actual collar slack
\[
 q=\frac1{2\pi}\int_{\alpha}^{2\alpha}P_V(t)\,dt>0
\]
and give budgets that keep the perturbation, tree error and transfer cost below
fixed fractions of $q$.  Independently choosing a shortest trajectory in each
strip is not enough unless the attachment mismatch is controlled.

\subsection*{4. Coefficient perturbation and stability}

The constant-translation stage is no longer open.  Once a finite critical-value
family is injective, Lean proves an arbitrarily small translation making every
value nonzero and pairwise positive-ray separated
(\lref{Erdos1041/NewtonFlowRaySeparation.lean}{197}{exists_small_translation_separating_arguments}).  It also proves the
explicit root-retention estimate; a shift below $\varepsilon$ keeps all roots
in the unit disc when
\[
 ((n+1)\varepsilon)^{1/n}+\rho<1
\]
(\lref{Erdos1041/NewtonFlowRaySeparation.lean}{287}{constant_perturbation_roots_in_unitDisk}).  A constant translation
cannot separate initially equal critical values.

\begin{problem}[two-stage generic perturbation with slack]\label{prob:perturb1041}
For fixed root discs, compact collar $K$, regular levels
$\alpha/2,\alpha,2\alpha,5\alpha/2$ and collar slack $q>0$, prove that an
arbitrarily small
\[
 g_{\lambda,\beta}(z)=f(z)+\lambda z+\beta
\]
can be chosen so that:

\begin{enumerate}
\item $f+\lambda z$ has simple relevant critical points and injective complex
critical values;
\item the checked constant translation $\beta$ makes those values nonzero and
pairwise ray-separated;
\item no critical point enters the protected collar or truncation boundary;
\item the relevant component remains in $K$, contains exactly the corresponding
perturbed roots and has slack $q_g>q/2$; and
\item root displacement and straight-line transfer back to the original roots
consume less than a prescribed fraction of $q/m$.
\end{enumerate}

If the one-coefficient perturbation $\lambda z$ cannot ensure all five
properties, give the smallest additional lower-coefficient direction that can,
or an explicit obstruction.
\end{problem}

Finite planar avoidance and the subsequent constant translation must not be
relisted as missing; coefficient genericity, component stability and slack
stability are the open content.

\subsection*{5. The global Newton-flow claim ceiling}

\begin{problem}[relative global Newton-flow theorem]\label{prob:globalflow1041}
On a compact regular band
\[
 M_{a,b}=\{z:a\le-\log|p(z)|\le b\},
\]
prove real-time existence and the identities
\[
 p(z(t))=e^{-(t-t_0)}p(z(t_0)),\qquad
 u(z(t))=u(z(t_0))+t-t_0
\]
throughout every maximal orbit; classify both limiting endpoints among the
regular boundary, root ends and finite saddle set; and exclude recurrent,
accumulating and non-Hausdorff orbit phenomena.  Under pairwise ray separation
of the critical values, decide whether the flow is Morse--Smale relative to the
boundary or whether the strongest valid conclusion is only absence of
saddle-to-saddle connections.
\end{problem}

The checked algebra proves the pointwise value equation and the endpoint-ray
consumer.  It does not supply global solution theory or an orbit-space graph.
A positive stronger theorem must give the graph and a finite edge bound; a
negative answer should exhibit the simplest ray-separated polynomial carrying
the remaining pathology.

Erd\H{o}s~\#1041 remains open.  Five separate statements are proved above,
with different degree ranges, root hypotheses, containment levels and
constants.  Theorem~\ref{res:low-critical-thirteen-twentyfifths} holds in
every degree $n\ge2$ for squarefree $f$ under the single hypothesis
$\mu\le13/25$, with no condition on where the roots lie, and gives a
connector inside $\{|f|<1\}$ of length below $2$; the regime
$13/25<\mu<1$ is untouched.  Theorem~\ref{res:constant-factor-path} drops
every threshold and every root hypothesis, and pays for that with the
constant $71/10$ and the weaker containment level $\{|f|\le2\mu\}$; the
containment $\{|f|<1\}$ with length at most $5.7$ needs roots in the open
unit disc and $\mu\le1/2$, and both fall short of the target constant $2$.
Theorem~\ref{res:degree-three} settles degree three completely for roots in
the open unit disc.  Theorem~\ref{res:critical-value-separation} holds in
every degree and is conditional on the critical-value separation $S$, with the
threshold $S=2$ reaching every degree $n\ge3$ by
Corollary~\ref{res:critical-value-thresholds}.
Theorem~\ref{res:trinomial-all-degree} settles the whole trinomial family in
every degree, and its radial inequalities are kernel-checked.  All five are
ordinary proofs at the level of the assembled path.  Their Lean companions check named numerical, algebraic and
power-series steps inside them, the degree-three companion lies outside the
pinned formal-source library and carries no kernel receipt in this release,
and the univalent branch, the area bounds and the certificate chain remain
ordinary mathematics.  The source now publicly verifies the Newton
kernel, finite ray avoidance and quantitative constant-translation root
control; the coefficient perturbation, corrected planar decomposition and
metric gluing remain the exact unresolved producers.

\section*{Statements and declarations}

Lean does not check the exposition, citation choices, or interpretation. This
manuscript cites Lean only for the formal statements and proofs that the pinned
kernel accepts. The checked
core is the Newton value equation, the exponential first integral, the consumer
form of ray separation, the finite planar-avoidance theorem, quantitative
constant-translation root retention, and the ray-collision parameterisation.  The
decomposition and length statements of \S\ref{sec:open} are not proved.  The
diagnosis in \S\ref{sec:gap} concerns the printed local saddle construction,
while the Cassini calculation refutes the displayed unrestricted tree-budget
estimate as an ordinary mathematical argument.  Attribution to a numbered
historical proposition remains conditional on its complete archived
hypotheses.  The
search results of \S\ref{sec:finite} are computations.

\appendix
\section{Guide to the formal sources}\label{app:sources}

The public \texttt{ErdosProblems.Erdos1041.NewtonFlowRaySeparation} module
contains the checked source for this note.  The search of \S\ref{sec:finite} is
\texttt{scripts/search\_counterexample.py} in the source package.  The
declaration table below is pinned to the shared formal-source commit used
throughout this problem-note series.

\begin{itemize}[leftmargin=*]
\item \lref{Erdos1041/NewtonFlowRaySeparation.lean}{34}{newtonFlowVector}
\item \lref{Erdos1041/NewtonFlowRaySeparation.lean}{38}{derivative_mul_newtonFlowVector}
\item \lref{Erdos1041/NewtonFlowRaySeparation.lean}{50}{newtonFlow_value_hasDerivAt}
\item \lref{Erdos1041/NewtonFlowRaySeparation.lean}{64}{newtonFlow_scaledValue_hasDerivAt_zero}
\item \lref{Erdos1041/NewtonFlowRaySeparation.lean}{77}{SamePositiveRay}
\item \lref{Erdos1041/NewtonFlowRaySeparation.lean}{80}{samePositiveRay_refl}
\item \lref{Erdos1041/NewtonFlowRaySeparation.lean}{84}{samePositiveRay_symm}
\item \lref{Erdos1041/NewtonFlowRaySeparation.lean}{92}{samePositiveRay_trans}
\item \lref{Erdos1041/NewtonFlowRaySeparation.lean}{107}{translated_samePositiveRay_parameterization}
\item \lref{Erdos1041/NewtonFlowRaySeparation.lean}{127}{realAffineLine}
\item \lref{Erdos1041/NewtonFlowRaySeparation.lean}{130}{realAffineLine_eq_affineSpan}
\item \lref{Erdos1041/NewtonFlowRaySeparation.lean}{147}{isClosed_realAffineLine}
\item \lref{Erdos1041/NewtonFlowRaySeparation.lean}{152}{dense_compl_realAffineLine}
\item \lref{Erdos1041/NewtonFlowRaySeparation.lean}{162}{exists_small_avoiding_finite_realAffineLines}
\item \lref{Erdos1041/NewtonFlowRaySeparation.lean}{179}{samePositiveRay_imp_mem_realAffineLine}
\item \lref{Erdos1041/NewtonFlowRaySeparation.lean}{197}{exists_small_translation_separating_arguments}
\item \lref{Erdos1041/NewtonFlowRaySeparation.lean}{230}{norm_lt_one_of_near_root}
\item \lref{Erdos1041/NewtonFlowRaySeparation.lean}{257}{constant_perturbation_root_near_original}
\item \lref{Erdos1041/NewtonFlowRaySeparation.lean}{287}{constant_perturbation_roots_in_unitDisk}
\item \lref{Erdos1041/NewtonFlowRaySeparation.lean}{325}{samePositiveRay_of_real_exp_decay}
\item \lref{Erdos1041/NewtonFlowRaySeparation.lean}{334}{no_newtonConnection_of_not_samePositiveRay}
\end{itemize}

\paragraph{Checked declarations behind the translated quotient-fibre families.}
For the translated cubic and quartic quotient families $f(z)=P((z-h)^q)$ the
metric theorems are ordinary; the kernel checks the cyclic-fibre mean square
and the quotient-disk bridge.  The fibre norm identity is
\lword{Erdos1041/CyclicFiberMeanSquare.lean}{21}{cyclicFiber_sum_normSq}{fibre
mean square}, the strict unit bound is
\lword{Erdos1041/CyclicFiberMeanSquare.lean}{37}{cyclicFiber_normSq_add_lt_one}{fibre
unit bound}, the quotient roots land in the open disk by
\lword{Erdos1041/CyclicFiberMeanSquare.lean}{60}{cyclicFiber_quotient_norm_lt_one}{quotient
disk bridge}, and the composition vanishes on the fibre by
\lword{Erdos1041/CyclicFiberMeanSquare.lean}{73}{cyclicFiber_composition_zero}{fibre
zero identity}.
